\documentclass[12pt]{book}

\usepackage{graphicx}
\usepackage{eso-pic}
\usepackage{tcolorbox}
\usepackage{polyglossia}
\setdefaultlanguage[indentfirst=false,forceheadingpunctuation=false]{russian}
\setotherlanguages{english}

\RequirePackage[a4paper, headsep = 0.5 \headsep, left=2.5cm, right=2.1cm, top=2cm, bottom=2.1cm]{geometry}

\setmainfont{URW Bookman}
\newfontfamily\cyrillicfont{URW Bookman}[Script=Cyrillic, Ligatures=TeX]
\setlength{\parindent}{20pt}

\usepackage{enumitem}
\setlist[itemize, 1]{label =\raisebox{-0.3\height}{\scalebox{1.8}{\textbullet}}}

\title{Феноменология чуждого}
\author{Бернхард Вальденфельс}
\date{2006}
\begin{document}

\AddToShipoutPictureBG*{\includegraphics[width=\paperwidth]{title.png}}

\null
    \thispagestyle{empty}%
    \addtocounter{page}{-1}%
    \newpage

\tableofcontents

\chapter{Грани чуждого}

Выделить чуждое как отдельную тему --- значит уже упустить его, начать со знакомого места и планировать к нему же вернуться. Обращение к чуждому наверняка приведет к переменам и даже катарсису, хотя в конечном счете привычность и возобладает и может даже углубиться. Чуждое не безобидно и может отдалить нас от самих себя --- потому ему стараются сопротивляться, избегать его или ассимилировать. Эти старания оставляют субъекта наедине с самим собой и требуют выстраивать и оборонять порядок, исключающий проявление чуждого изнутри. Чуждое может пробудить любопытство и рассказать нам о нас самих, но как только оно подходит к тайне свободы и разума --- срабатывает сигнал ''Хаос!'', и свобода и разум остаются в своих правах. Иначе им пришлось бы отказаться от самих себя. Чуждость неизбежно ведет к враждебности, обе стороны все больше убеждаются в своей правоте. Мы становимся ближе и бдительней друг к другу. Встроенные в эмпирическую сеть защитные механизмы исходят из \textit{своего} и стремятся закрепиться в \textit{общем}. Чуждое предполагается внешним, ожидается, что оно будет идентифицировать себя и выглядеть незваным гостем. Его тогда можно будет подвергнуть оценке и суждению и поместить на повседневный моральный, политический, религиозный, культурный и интеллектуальный карантин.

Если же чуждое воспринимать как неизбежно неопределенное, как всегда сокрытое в \textit{своем}, где оно тревожит, соблазняет или пугает, ускользает от понимания, то опыт чуждого, значит, воздействует на \textit{наш} опыт, становясь отчуждением от опыта. Чуждость самореферентна и заразительна, ее последствия предшествуют всякой ее тематизации. Феноменология чуждого, предвосхищенная Гуссерлем в нескольких его новаторских концепциях (чуждый опыт, чуждое эго, чуждое тело и чуждый мир) обладает задатками науки, могущей говорить о чуждом во всей его полноте, всерьез принимать его радикальность, не довольствуясь сведением к его пониманию и практическому распознаванию. Феноменолог, всерьез принимающий эти требования ксенологии, наверняка будет вдохновляться работами Марселя Мосса, Георга Зиммеля, Вальтера Беньямина, Людвига Витгенштейна или Михаила Бахтина, не говоря уже о таких литераторах как Кафка, Музиль, Целан, Валери и Кальвино. В мире искусства другие порядки, и художественное чутье легко преодолевает границы самоуверенности и принципиальности, если вступает в контакт с чуждым.

Настоящее исследование не пытается охватить проблематику чуждого во всей ее пугающей полноте, а лишь обозначает некоторые основные мотивы, побудившие автора разработать феноменологию чуждого в соответствующей серии книг. Первый том --- \textit{Жало чуждого} вышел в 1990г.\footnote{Формальное изложение феноменологии чуждого можно найти в Главе 1 \textit{Топографии чуждого} (\textit{Topographie des Fremden}, 1997). Аспекты чуждого рассматриваются в главах 5-6 книги \textit{Рваные линии опыта} (\textit{Bruchlinien der Erfahrung}, 2002).} Мотивы здесь развивают тему, а не просто ретроактивно обосновывают обращение к ней. Опыт чуждого --- как и опыт удивления, ужаса, любви, ненависти или тех едва заметных несостыковок, через которые возникнет новое --- нельзя вполне описать или обосновать. Разделение на главы выполнено так, что каждая рассматривает конкретный программный компонент проблематики чуждого. Они выстроены в порядке этапов формирования феноменологии чуждого --- они потому вводные, поясняющие тексты и дополняют уже опубликованные обширные исследования. В последующих абзацах я свяжу эту работу с предыдущими, что избавит меня от необходимости отдельно на них ссылаться и позволит сосредоточиться на ключевых темах: порядок -- пафос -- ответ -- тело -- внимание -- межкультурность.

В главе 1 рассматривается \textit{Порядок в сумерках} (\textit{Ordnung im Zwielicht}, 1987). На границах всякого порядка чуждое предстает экстраординарным, не находящим места в порядке --- и все же оно, будучи исключаемым, не ничто. Поскольку оно исключается не вообще --- а из конкретного порядка, то оно нечто большее, чем просто сливающееся с фоном в неопределенности. И поскольку оно взывает к нашим чувствам, то оно не просто шум, от которого мы хотели бы укрыться. Положение, в котором мы оказываемся, обнаруживая несхватываемое в схватываемом, беспорядочное в упорядоченном, невидимое в видимом, тишину в озвучиваемом, запускает цепь исследований, сводящихся к вопросу: как мы можем взаимодействовать с чуждым без предварительной нейтрализации или отрицания его эффектов, его вызовов и требований? Как феноменолог, я всегда старался избегать того, в чем нетерпеливая аудитория любит обвинять феноменологов чуждого --- возведения его в ранг некоего морального или религиозного авторитета или ''совершенно иного''. Это было бы обожествлением Другого --- в чем часто обвиняют Левинаса --- и морализацией чуждого. Но навязать переживание столь же возможно, как и, согласно Канту, чувство. Не то чтобы восприятие чуждого было лишь субъективным ощущением, как в том убеждает модерн, просто в этом мире вещей гораздо больше, чем может охватить школа мысли с ее нормами, ценностями и интересами.

К переживаниям можно относиться с разной степенью осторожности: позволить им говорить самим за себя, использовать как информационный ресурс или пытаться исключить --- будь то от небрежности, властолюбия или безразличия. В главах 2 и 3 эти возможности рассмотрены с точки зрения ответностной этики --- этики, возникающей в ответе, отклике, реакции, и так остающейся за рамками заповедей и запретов. До-предикатный опыт в гуссерлевском смысле здесь обращается в до-нормативный --- остающийся всегда позади, пока мы намечаем действия и строим планы, требующие моральной или правовой оценки. В главе 2 я буду ссылаться в основном на \textit{Рваные линии опыта} (\textit{Bruchlinien der Erfahrung}, 2002): анализ, проведенный в этом тексте, построен вокруг двойного события, состоящего из пафоса и последующего ответа. В начале есть то, что предстает как события --- происходящим с нами, в чем я или Другой можем участвовать --- а не как волящий и действующий субъект. Ключевое здесь --- временной сдвиг. Не нечто иное предшествует инициативе ''я'' --- эта инициатива предшествует самой себе, не возникая из неоткуда. Эти связи предстанут в новом свете, когда в главе 3, рассматривая \textit{Регистр ответов} (\textit{Antwortregister}, 1994)\footnote{В данном случае я намеренно отклоняюсь от порядка, в котором были изложены эти идеи, чтобы облегчить понимания, поскольку \textit{Регистр ответов} предполагает некоторую часть того, что полностью раскрывается только в \textit{Рваных линиях опыта}. ''На что'' ответа указывает на ''чем'' преодоления.}, я раскрываю основные черты ответового поведения и возникающую в результате зова чуждого логику ответов. Парадокс творческого ответа представляет особый случай: оказывается, мы можем воспринимать ''что'' ответа как нечто изобретенное, но не можем сделать то же с его ''на что''. Зов чуждого может быть столь же слабо связан с этическим фундаментализмом, сколь и с технологическим конструктивизмом.

Глава 4 посвящена телесному опыту, и не в коем случае не просто расширяет сферу чуждого. В узле самореференции и самоустранения, определяющих его способом существования, тело оказывается как бы эмблемой чуждого. Чужеродность начинается не с меня самого --- это мало что значило бы,--- так называемый субъект выпадает из центра, когда предстает как \textit{Воплощенное Я} (\textit{Leibliches Selbst})\footnote{Название лекций о теле, опубликованых в 2000г., в которых эта тема рассматривается во всей ее глубине.}. В главе 5 уточняются и подкрепляются ссылками на литературу основные идеи \textit{Феноменологии внимания} (\textit{Phänomenologie der Aufmerksamkeit}, 2004). Внимание полностью подчинено двойному ритму становления очевидным и становления притягательным, двойное движению пафоса и ответа потому преобретает форму, пронизывающую весь наш опыт и поведение. Что становится очевидным или мыслимым, что доходит до наших глаз или ушей, до наших органов чувств, что до ходит до нас --- не наше. Аллюзия названия \textit{Пороги смысла} (Sinnesschwellen, 1999) указывает, что вездесущее внимание ведет в область эстетики и особого этоса чувств. В последних двух главах говорится также о чуждых аспектах психоанализа и технологии, что занимало значительное место в \textit{Рваных линиях опыта} и \textit{Феноменологии внимания}. Как изначальный опыт пронизан моментами неизведанного и немотивированного, так и феномены в нем с самого начала отмечены подлинной феноменотехнологией и соматотехникой. Технология очень долго --- вопреки всем сетованиям современности об отчуждении в ней --- рассматривалась как расширение сферы собственности, но теперь все больше становится источником особой отчужденности.

Главу 6 можно читать как обзор, собрание различных аспектов чуждого в кратком, ориентированном на отношения между культурами, изложении. Межкультурность не ограничена философией.\footnote{Позвольте здесь упомянуть о выходящем в Вене журнале по межкультурной философии Polylog, в котором к настоящему времени появилось довольно много тематических выпусков.} Глобальное мышление не следует ни ожидать, ни желать --- оно неизбежно сформировало бы лишь новые гегемонистские отношения, разве что другими средствами. Философия, переставшая вращаться вокруг себя самой и пустившая корни в привычной культуре, дала бы облегчение. Попытка пересечь границы, не уничтожив их --- одно из будущих начинаний в области межкультурного взаимодействия, и оно значит нечто большее, чем просто сопоставление или смешение культур. Разнообразие слишком легко становится накоплением --- спокойным перебором любопытных вариаций. Опасность для культур, достигних зрелости, заключается в обращении к обществу охотников и собирателей.

Настоящая работа основана на серии отдельных исследований, появившихся в связи с разными обстоятельствами. Все эти исследования были переработаны и ориентированы на общую тему.\footnote{Эти идеи рассматриваются в \textit{Границы нормализации} (\textit{Grenzen der Normalisierung}, 1998).} Я постарался избежать ретроактивной систематизации и сохранить эвристическую множественность подходов, в то же время создав рельефы, благодаря которым --- как на объемной открытке --- были бы видны все значимые аспекты. Как писатели, мы иногда сожалеем, что письмо ограничено двухмерной плоскостью, но, к счастью, в письменной речи --- как и в чтении --- есть свои способы расставления акцентов, и эти подъемы и спуски не дают определенному направлению письма свестись к написанному.

\chapter{Человек как пограничное существо}

Чуждое --- это предельное как оно есть. Даже возникая в знакомом, оно приходит извне. Не бывает чуждого без чуждого места, и порядок нашей жизни, опыта, языка и действий определяет, сколь значимо для нас будет чуждое. Преобразование порядка значит и преобразование чуждого --- такого же многогранного, как и соответствующий порядок. О чуждом говорят не меньше, чем об эго --- пограничные зоны между и за порядками питают чуждое.

\section{Порядки и их пределы}

Обычно верно полагать, что не только порядок ограничен, но и что эти границы возникают в процессе упорядочивания. Нечто --- то, что оно есть в силу его отделения от прочих вещей: камней, растений, животных, природных или искусственных объектов. Однако, в мире людей понятие границ вызывает особое беспокойство, потому что границы всегда под сомнением. Поведение человеческого существа не определяется ни инстинктами, ни искусственными программами: оно не ограничено, а, скорее, особым образом соотнесено с границами. Это касается и места и времени --- границ, определяющих что такое здесь и сейчас,--- и запретов, сдерживающих желания и поступки, и пределов понимания, сдерживающих мышление. Неудивительно, что вопрос о границах существования остается в числе главных тем на протяжении всей человеческой истории: Иегова отделяет свет от тьмы, границы души неизмеримы, философ намечает границы чистого разума, системный теоретик воплощает возвышенный жест созидания в простом ''Проведи границу!''.

Видно, что установление границ, приводя к разным порядкам и структурам, носит не только прагматический и региональный, но и эпохальный характер. Каждая эпоха --- или каждая культура, общество, среда или форма жизни --- ведет себя в рамках некоторых границ, и отношение к этим границам, всегда сопровождаясь определенной политикой, значительно варьируется. Обращение с границами прямо показывает глубинный дух эпохи и может также послужить комментарием к тому, что так долго продвигало современность вперед или предшествовало ей, или подрывало ее и выходило за ее пределы. Рубеж к 21-ому столетию --- определенно очень специфическая граница, и у нас нет языка, подходящего для описания того, что ждет нас впереди.

\section{Безграничная вселенная}

Для контраста можно начать с безграничной вселенной. В нашей культурной традиции она наиболее емко представлена греческим космосом. Космос долгое время имел парадигмальное значение, олицетворяя классическую форму порядка. Он воплощал не один из множества порядков, а порядок вообще, как таковой, и потому единственной альтернативой ему были неупорядоченное многообразие хаоса. Этот космос придает каждому существу его собственную форму (πέρας), границы которой отделяют существо от себя и среды. Точная форма выражается в концептуальном определении (όρισμός); со времен Платона диалектика связывала каждое бытие как единое с прочим в цепочке отношений. Эта горизонтальная связь дополняется вертикальной иерархией, пропорции которой соответствуют степени отражения рационального целого в конкретном существе. Люди здесь стоят выше животных, греки --- выше варваров, мужчина --- выше женщины, созерцание --- выше действия. Рациональность, раскрывая закон целого, имеет решающее значение для места человека в иерархии существ.

Эта реляционная структура --- все различия в ней относительны --- имеет нижнюю и верхнюю границы. Нижняя формируется \textit{ин-дивидом}\footnote{in-dividual; англ. divide --- разделять.}, άτομον είδος --- он уже не может быть далее разделен без разрушения его идентичности. Верхняя граница --- Вселенная, не упорядоченная в рамках чего-либо другого; мир --- сущность, из которой нечего не исключено, постановит Аристотель в своей \textit{Физике} (2.6.207a8). Проще говоря, космос --- порядок без внешнего, с только лишь внутренними границами. И пересекающий эти границы либо попадает в дурную бесконечность бездонного и бессмысленного апейрона, либо, подобно Икару, поднявшемуся на запредельные высоты, терпит фатальную неудачу.

Но замкнутый всеобъемлющий космос, в котором есть место каждому, предполагает, что место проявления и выражения целого как \textit{целого} находится внутри самого этого целого. Психе --- согласно Аристотелю, ''в некотором смысле --- все'' --- становится ступенью порядка, к которому мимически приспосабливается. Космос потому предстает порядком, раскрывающим и выражающим \textit{себя как себя} через превращение всякого условия в момент себя. Такой порядок без внешнего соответствует \textit{мысли о внутреннем} --- \textit{penser du dedans}, перефразируя название работы Фуко\footnote{La Pensée du dehors, 1966} --- и в целом такая мысль была бы мыслью самой по себе. Однако даже древние греки знали о выпадающих из этой схемы фигурах: Сократе --- άτοπος, живущем в городе подобно прочим гражданам, но действующем как вопросительный знак, или о Платоновской \textit{мании} поэзии, или о софистах, выступающих против истинного логоса лексическими уловками, или об Эдипе, прозревшем только став слепым, прикованном к своей зловещей судьбе, изолировавшей его и заклеймившей бездомным (άπολις).\footnote{Фигура \textit{аполиса} вырисовывается на фоне полиса, и Софокл сравнивает его с ύψίπολις, ''самым выдающимся в городе'' (\textit{Антигона} 5.370). Человек-\textit{аполис} соответствует ''лицу без гражданства'' в нынешнем понимании, но в нем также виден аналог \textit{атопоса}: ''лишенный места'', ''неуместный'', ''странный''. Alpha privativum в обоих случаях можно понимать как не только отрицание, но и выражение отказа, отстранения, выхода.} Такие маргинальные фигуры встречаются и в других местах: приверженность текстам в общепринятой части мистического учения обычно сопровождается ответвлениями, отбрасывающими тень неортодоксальности и анархии --- это касается и еврейской и исламской, и христианской традиций. Даже Святой Павел отмежевывается от риторики, в которой γλώσσα или \textit{lingua} расслабляют язык до такой степени, что позволяют ему ускользнуть от индивидуального или общественного контроля. Высказывание, растворяющееся в космосе сказанного --- заменяется ничего не говорящим и потому бессмысленным. Обаяние, противопоставляемое Максом Вебером повседневной жизни институтов, пронизывает все и показывает, что любая нормальность --- даже космологическая, теологическая или космополитически моделируемая --- упускает нечто, находящее затем выражение в некой безумной бахроме аномалий.

\section{Принадлежность, чуждость, контингентность}

Глобальный порядок немыслим без разрыва, если только место, из которого разворачивается целое, не превращается в конкретное целое. Модерн --- движение, ставящее под сомнение видение целого. Оно пресекает самообман космотеороса, считающего себя частью созерцаемой им сцены, и мнимый порядок оказывается результатом процесса установления порядка [\textit{Ordnungsstiftung}]. Здесь два ключевых открытия: (1) \textit{я}, говорящее ''я'' до того, как будет названо ''субъект'' и взрыв структуры отношений в целом из-за этой самореферентности и (2) радикальная \textit{контингентность}, которая не только занимает открытые пространства в порядке, но и влияет на сам порядок. Такой порядок может не только выродиться в беспорядок, но и стать другим порядком, быть отличным от самого себя. В cogito Декарта есть не только вершина мысли, но и возможность создания Богом другой математики. Порядки, в которых мы движемся, оказываются лишь потенциальными. И странный протагонист \textit{Человека без свойств} Музиля подозревает, что даже Бог предпочитает говорить о мире лишь потенциально: ''ибо Бог создает мир и в то же время думает, что он мог бы быть другим'' (Музиль, 1961, 15) и связывает это с тезисом ''действительно любящий свою страну не считает ее лучшей''. Решающий фактор здесь --- что мотивы субъектности и рациональности выступают как двойной мотив. Субъект модерна больше не олицетворяет лишь рациональность и ищет свое место, которого нет.

Если всерьез --- во всей его радикальности --- принять этот двойной мотив, то проблема проведения границ изменится: \textit{я}, живущее в своей собственной сфере и своему собственному культурному порядку соответствующее даже в языке и чувствах, уже не будет сводимо к чему-то всегда одному и отдельному от другого в контексте универсальности. Самость и самопринадлежност --- результат проведения границ, разделяющих внутреннее от внешнего и так принимающих формы \textit{включения} и \textit{исключения}. Чувство принадлежности возникает, когда что-то выходит из него, и это уходящее мы принимаем чуждым и неоднородным. Такое разделение своего и чуждого осуществляется без участия третьей стороны и составляет измерение, отдельное от ''то же / другое'' (которое опирается на диалектически созданное целое). Выражаясь языком западной философии, другой (έτερον) и чуждый (ξένον) --- не одно и то же. Чуждость гостя --- в том числе Незнакомца из Элеи, появляющегося в \textit{Софисте} Платона,--- иного языка или культуры не сводится к тому что нечто кажется другим. Дерево и бетон отличаются друг от друга, но никто не станет говорить, что они чужды друг другу. Чуждость значит, что у \textit{я} (\textit{ipse}) есть сфера принадлежности, собственного бытия, и что это \textit{я} не то же что \textit{то же} (\textit{idem}), которое различимо лишь третьей стороной.

\section{Компромиссы модерна}

Ницше пишет: ''Важно избавиться от Вселенной, от единства ... разрушить Вселенную, разучиться уважать ее'' (KSA 12:317). И едва ли модерн серьезно относится к последствиям разграничения порядков. До сих пор мы стремимся к компромиссам через взаимное смешение \textit{я} и \textit{того же} и \textit{другого} и \textit{чуждого}. Двойственность оценки потенциала нового порядка видна и в конфликте модерна и постмодерна. Взять хотя бы \textit{я}, которое говорит ''я'' и так отделяет себя от целого. ''Я'' буквально значит исключение. Прямое соединение картезианского эго с эгоцентричностью было бы чрезмерным упрощением, ведь пренебрегло бы беспокойством, проистекающим из этого радикального самоанализа и побуждающим философа искать незыблемый фундамент. Поиск обречен на провал, но что им движет? Даже если ''\textit{Кто} я?'' слишком быстро соскальзывает к ''\textit{Что} я?'', этот вопрос, замыкая вопрошающего и вопрошаемое, уже содержит след ''Je est un autre'' Артюра Рембо.

Как бы то ни было, открытие \textit{я}, чья самость отделяет его от структуры отношений в природном и социальном мире, ослабевает, когда его превращают в лишь \textit{частность}. В конце концов мы бы столкнулись только с конфликтом \textit{индивидуализма} и \textit{холизма}, \textit{партикуляризма} и \textit{универсализма} --- конфликта все еще сохраняющегося но никогда не приводившего к значимым результатам. Если целое нельзя понять без частей, в которых оно выражено, а общие правила малопригодны без учета обстоятельств частностей, к которым применяются, то конфликт этот приводит к великому объединению: аристотеликов --- с кантианцами, герменевтов --- с универсальными прагматиками. Но гармония нарушается, если вспомнить, что элементы ''Я-Ты-Здесь-Сейчас'' не принадлежат к общим концептуальным категориям --- скорее, эти индексальные или окказиональные выражения являются указателями на место произнесения речи, место развертывания поля опыта, языка и действия, прежде чем оно само займет какое-либо место. Место, о котором говорится, не совпадает с местом произнесения речи, и то же касается момента времени. Даже как некто, говорящий ''я'', я не являюсь счетным элементом класса или частью целого. ''Я'' произношения, будучи явно отлично от ''я'' содержания произнесенного --- не исчислимая вещь, о которой можно было бы говорить во множественном числе, и то же касается ''ты''. В древнегреческой мысли ''я'' и ''мы'' не играют тематической роли за исключением понятия ''для нас'' (πρός ήμας). Взрывная сила я-речи также теряется, если сводить ''я'' к общей я-функции. Здесь нам нужно защитить Декарта от его трансцендентально-филосовских наследников.

То же затруднение касается и второго --- приписывающего всем порядкам неустранимую контингентность --- открытия. Если рассматривать \textit{способность быть другим} или \textit{способность действовать по-другому} как большую или меньшую степень \textit{произвольности}, то антиподом полагается \textit{необходимость}. Касательно претензий на легитимность, это ведет к еще одному конфликту --- релятивизма и универсализма, оружие которых сейчас уже столь же изношено. Здесь стоит еще раз напомнить: установления порядков с их легитимностью --- включая генеалогию истинного и ложного или добра и зла --- не верно ни относительно, ни абсолютно. Оно никак не верно, потому что существование бинарных стандартов не подчинено самим этим стандартам, если только их генезис опять не скрыт и соответствующая оппозиция не гипостазирована. У каждого порядка есть свое слепое пятно --- что-то неупорядоченное, что не просто нехватка. Это относится как к моральным, так и к когнитивным и эстетическим порядкам. Поэтому современные авторы так упорно борются с \textit{инципитами} своих романов: с самого начала и писатель, и читатель попадают в ловушку готового порядка.\footnote{См. \textit{Феноменология внимания} (\textit{Phänomenologie der Aufmerksamkeit}), глава 3 \textit{Неисчислимое} (\textit{Unerzählbares}).} Этот порядок постанавляет: ''есть порядки'', где ''есть'' не поддается никакому обоснованию, являясь базовым допущением всякой попытки обоснования --- нулевым уровнем. То, что сейчас называется постмодернизмом, может значить, среди прочего, что некоторые поля сражений утратили свое значение --- но не гарантирует, что проблемы, проистекающие из модерновой субъектности и рациональности, не возникнут вновь уже в другом месте.

\section{Парадоксы самоограничения}

Парадоксы нынче в моде. Парадокс --- παρά δόξαν --- буквально мнение, расходящиеся с преобладающим. Как и в случае других пара- (паразитизм, парафизика, паралогизм, парапсихология), понятие парадокса предполагает нормальное положение вещей. Покуда парадоксы отклоняются лишь от преобладающих мнений, к ним нет возражений. Совсем другое дело, когда предположение вступает в противоречие с собственными предпосылками, усиливая одно предположение за счет другого, и наоборот. Для сохранения предположений значение парадоксов обычно преуменьшают, как в случае с парадоксом лживого критянина. Конечно, различение объектного и метаязыка спасает высказывание --- его событие теперь не часть его содержания. Но такое методологическое решение основано на вышеупомянутом ограничении, приходящим извне и требующем различения между двумя уровнями речи. Анализ и решение проблемы находятся на уровне говоримого, что значит пренебрежение возможностью высказывания и само-высказывания --- сказанного, но не явленного в речи. Самореферентность есть, например, в ''Я обещаю вам, что ...'', из-за нее критянин вовлекает себя во все, что говорит. Кроме того, хорошо известное различие между плоскостью содержания и плоскостью отношений из теории коммуникации Пола Вацлавика не может быть прямо применено \textit{к} речи, потому что всегда присутствует \textit{в} ней.

Нечто подобное происходит и при самоограничении, когда ограничивающее себя является результатом ограничения, а не некой третьей стороной, отделяющей одно от другого, не будучи в них вовлеченной. Эксперт по грибам ни ядовит, ни съедобен, и судья не имеет отношения к деяниям обвиняемого и ущербу истца. Но что происходит, когда некто извиняется? Легко вспомнить случаи, когда мы сталкивались с чем-то странным, от чего \textit{сами} ограждаются --- без третьей стороны: разница между близостью и удаленностью как разница между ''здесь'' и ''тогда'', ''сейчас'' и ''давным-давно'', чередующиеся миры бодрствования и сна, сверхъестественные ''беседы с умершими'', переносящие из одного мира в другой через провал в памяти в романе Хуана Рульфо ''Педро Парамо''\footnote{Интересный взгляд на мотив чуждого в современной литературе и теории литературы см. в \textit{Мексика за гранью одиночества --- попытка межкультурного анализа} (\textit{Mexiko jenseits der Einsamkeit -- Versuch einer interkulturellen Analyse}, 1994).}, отношения между мужчиной и женщиной, ребенком и взрослым, социальная изоляция из-за классовой принадлежности, профессии или культуры, наконец, границы нормы, отделяющие здоровых от больных, невиновных от преступников, ортодоксов от еретиков, своих от чужих. Во всех этих случаях относящий себя к одной из сторон оказывается вовне. Такие перемены, как засыпание и пробуждение, болезнь и выздоровление, старость, обращение в другую веру значат не новый вид от третьего лица, а становление чем-то другим. Пруст показал, в какой степени пробуждение зависит от чисто телесной самоориентации, его \textit{Исследования} начинаются на грани сна и бодрствования. Такие переживания --- когда свое сталкивается с чуждым --- преимущественно переживания пересечения границ и, как станет видно, сопровождаются удвоением движения. Двойная природа ограничения, когда, с одной стороны, некто отделен от другого, а с другой --- определенные области разграничены, допускает дальнейшую дифференциацию.

Первый аспект касается самореферентности действия по ограничению. Действующая граница не просто находится или проводится. Иллюстративно проведенная на доске линия здесь вводит в заблуждения. Акт проведения границы, когда некто отделяет себя от другого, нельзя ни увидить, ни потрогать --- он схватывается только через свой след. Его можно сравнить с заключением контракта --- действием, не становящимся частью самого контракта, но являемым косвенно, в изменении обязанностей. Процесс проведения границы происходит в нулевой точке, не находящейся ни здесь, ни там. Действующая граница потому --- ни определенная вещь, ни ничто, ведь без нее не было бы ни того, ни другого --- ни меня, ни других. \textit{Самореферентность проведения границ заключается в их самоустранении.}

Второй аспект касается самореферентности, проявляющейся в \textit{само} самоограничения и приводящей к внутреннему и внешнему ограничению. Эта самость --- ни завуалированный субъект собственного акта ограничения, ни объективный результат чуждого акта ограничения: она как бы выскакивает в самом акте ограничения, проявляясь как полость, \textit{как внутреннее, отделяющее себя от внешнего} и так создающее \textit{предпочтение в этой разнице}. Формально это значит: то, \textit{что} отделяет себя, отмечается, тогда как то, \textit{от чего} оно отделяется, остается неотмеченным. В этом неустранимая асимметрия, без которой не было бы себя, способного принять точку зрения другого или третьей стороны. И вышеупомянутые примеры иллюстрируют эту --- присущую контрасту между своим и чуждым --- односторонность. ''Здесь'' и ''сейчас'', отделяя себя от расстояния, прошлого и будущего, относятся к разграничению, но они же представляют место, в котором это разграничение происходит. То же верно для всякого другого разграничения. Мы не имеем дело с нейтральной, бестелесной, бесполой и нестареющей третьей стороной, которая из своей нейтральности проводила бы различие между мужчиной и женщиной, взрослым и ребенком --- нет и транснациональной или транскультурной третьей стороны, могущей, например, настроить немцев против французов или европейцев против азиатов. Немцы, отграничивая себя от французов, могут быть теми, кто они есть, только благодаря этому разграничению. Выделяющий, обособляющий себя --- на одной стороне, чуждое, непонятное, от чего он обособляется --- на другой. \textit{Обращение к чуждому заключается в отстранении от него.} И удваиваясь и множась, эта асимметрия не приводит к симметрии. Симметрия требует подчиненности своего и чуждого общей точке зрения или общему правилу. ''Уравнивание неравных'', упомянутое Ницше в ''Об истине и лжи во вненравственном смысле'' предполагает равенство \textit{не}равного. Нормализация как избирательная форма упорядочения, начинаясь уже с ортоэстезии органов чувств, не может быть понята без влияния гетероэстезии, и то же относится к ортодоксии, ортологии и ортопраксии.\footnote{Эти идеи рассматриваются в моей работе \textit{Grenzen der Normalisierung} (\textit{Пределы нормализации}) (1998)}

Парадокс самореферентности усиливается еще и тем, что первый аспект самореферентного ограничения и второй --- заключающийся в ограничении себя внутрь и так ограничении чуждого вовне --- перетекают друг в друга в форме \textit{обращения к чуждому в обращении к себе}. От неизбежной самотематизации, ''я'' оказывается в системе упорядочивания, которую само же создает. ''Здесь'' тела можно определить, например, по красному маркеру на карте, указывающему местоположение пользователя. ''Сейчас'' имеет дату и отмечено в календаре. Тело (будет подробно рассмотренно в следующей главе) проявляет себя как живое тело, постоянно ускользающее от рефлексии сознания и сопротивляющееся нашим инициативам. То же происходит на социальном уровне и проявляется в межкорпоральности. ''Я'', без которого не было бы других, описывается Гуссерлем в \textit{Кризисе} как ''привелигированный член'' (\textit{Hua} VI:188; \textit{Crisis}, 185), но в то же время ''я'' обнаруживает себя обычным членом группы, которую помогает конституировать. ''Я'' --- это Другой, потому что чуждость рождается в своем. Обращение к чуждому в обращении к себе объясняет, почему никто не является просто тем, кто он есть, цепочка удвоения себя происходит при различных обстоятельствах в работах Гуссерля и Мерло-Понти, Фуко и Лумана. Эти удвоения не следует путать с отражениями рефлексивного самосознания, являющегося и субъектом, и объектом, не связаны они и с интерсубъективной диалектикой узнавания --- когда один субъект обнаруживает себя в другом субъекте. Телесное самоустранение, становясь очевидным в различии между функционированием и тематизацией, соответствует скорее различию между актом выражения и выраженным, между высказыванием и сказанным, ведь событие высказывания никогда не исчерпывается сказанным. Вообще, такого рода беспокойство, приводящее к постоянному удвоению ''я'' проистекает из того, что место, откуда проводятся границы, находится ни внутри, ни за пределами порядков, но сразу и внутри и снаружи них. Самоустранение значит опасность определенных моментов чужеродности в ''я'', в данном порядке. Двойная игра, как показывает Фуко, превращается в обманку, когда мы пытаемся уловить видение в виденном, высказывание в сказанном, мышление в мысли, вместо восприятия этих удвоений как части самой игры.

\section{Возможности и невозможности}

Тип мышления, подчиняющийся парадоксам самоограничивающегося --- и, значит, самопревосходящего --- порядка, порождает ряд моделей мышления, не вписывающихся в рамки традиционного мышления. \textit{Отклонения} от порядка играют важную роль не только в современной поэтике, но и в самоорганизации природы, они не указывают на существование нечта отклоняющегося, потому что отклонения проистекают скорее из самоизменения зрения или речи, которые ускользают от себя. Другой тип мышления --- мышление \textit{избытка}. Этот тип речи и действия, двигаясь вдоль границ их соответствующих порядков, отстает от самого себя, потому что привносит в игру больше возможностей, чем может использовать, и в то же время выходит за пределы себя, потому что касается невозможностей невидимого, неслышимого и немыслимого. Большее или меньшее количество смысла относится к формам неклассического искусства и указывает на избыток или нехватку смысла, чего не найти в классических паттернах исполнения желаний и приказов. Как ''не-фиксированное животное'', вынужденное изобретать и создавать порядки, человек раскрывает себя и как испытывающего нехватку, и как обладающего избытком: не может положиться даже на естественные потребности, не может сосредоточиться даже на требующих единого стандарта, регулятивных идеях. ''Есть порядки'' относится к тем порядкам, которые действительно подчиняются необходимым и ограничительным условиям, для которых, однако, не найти достаточных оснований. Порядки создают возможности и невозможности, но возможным не становится установление новых порядков. Здесь находится момент фактической безусловности среди опыта. Радикально чуждо то, что невозможно предвидеть ни субъективными ожиданиями, ни транссубъективными условиями возможности. Избыток опыта, выводя нас за рамки существующих порядков, в конечном счете приводит к пространственно-временному \textit{сдвигу}: самореференция никогда не возвращается к себе, место и время речи никогда не совпадает местом и временем \textit{в} речи. Самоустранение значит присутствие отсутствием, близость отдалением. Сдвиг начинается с рождения, рождение никогда не полностью собственное, ведь никогда не переживается активно, не выбирается. Причастность к собственному прошлому игриво выражает Лоренс Стерн, когда раз за разом откладывает рождение своего героя Тристрама Шенди, в то же время предвосхищая его во фразах вроде ''Я еще не родился'' (1950, стр.33). Карло Леви пишет об этом: ''Тристрам Шенди не хочет рождаться, потому что не хочет умирать'' (Italo Calvino, \textit{Lezioni Americane}, 1988, 46). Рождение касается меня, и все же не мой акт. Оно всегда преждевременно, новорожденный всегда запаздывает, задержка эта присуща всему новому, всему выходящему за рамки существующих мер. Каждое новое формирование проявляется как \textit{реформирование} существующих формаций, потому первая речь и первый поступок невозможны --- как и последняя речь и последний поступок. ''Абсолютное настоящее'', которое заключало бы в себе весь смысл --- один из фантазмов традиционных порядков, отрицающих свое происхождение.

Чуждое, выходящее за пределы игрового поля возможностей как нечто экстра-ординарное, можно потому рассматривать как нечто не-возможное --- невозможность здесь не онтологическая, эпистемологическая, практическая или логическая, а проживаемая.\footnote{Эта невозможность не ограничивается патологической стороной, а включает ее. Мой вклад см. в \textit{Gelebte Unmöglichkeit} под редакцией C.Rohde-Dachser и F.Wellendorf (2004) в \textit{Inszenierungen des Unmöglichen: Theorie und Therapie schwerer Persönlichkeitsstörungen}.} Дефис значит, что выходящее за рамки предписаний ведет не в потусторонний мир, а за пределы этого мира. Однако вопрос о возможностях и невозможностях опыта, речи и поступков демонстрирует двусмысленность, с которой наши рассуждения достигают своей крайней точки. В связи с этим возникает вопрос о связи между самоосознанием и обращением к чуждому. Это обращение можно истолковать двояко: как \textit{ограничение} своих возможностей или как \textit{сомнение в своем}.

Герменевты вроде Гадамера или прагматики вроде Рорти, полагаясь на принадлежность к традиции, занимают неопределенную промежуточную позицию: чуждое включается в обсуждение ровно настолько, насколько оно ему принадлежит, так что любое сомнение в своем обращается сомнением своего в своем, каждый остается сам по себе. И слабость рассудочной способности традиции, о которой говорит Джанни Ваттимо, мало что меняет. С другими теоретиками все однозначней. Первую интерпретацию обращению к чуждому мы находим в системной теории Никласа Лумана, он настойчив и в то же время гибок в размышлениях о том, как обращаться с границами, но всегда останавливается на определенной границе. Для Лумана обращение к чуждому значит претерпевшее сдвиг обращение к себе. Каждое наблюдение, что-то знача, и тем отделяя его от прочего, само по себе может стать предметом наблюдения для себя или другого. Единственное странное, что остается --- operatio pura наблюдателя, смутно напоминающее actus purus из аристотелевской теологии. Слепое пятно наблюдения не стирается, а перемещается с одной системной позиции на другую. В отличие от коммуникативного разума, всегда превращающего странное в нечто общее, системный оставляет место для чуждости, но чуждость эта --- \textit{исключенные, но при определенных условиях включаемые вновь, возможности}. Она лишь обозначает границы способностей --- личных, общих или оперативных. Таким образом, самореферентность первична, даже когда доступ к себе четко ограничен, и эта первичность превращается в свою полную противоположность, когда ''я'' подчиняется силе чуждого и так признает его своим безусловным хозяином. Похоже, обе крайности --- \textit{функционалистское непонимание} чуждого и \textit{фундаменталистская идеализация} чуждого --- стремятся усилить друг друга. Тогда мимолетной игре с возможностями была бы противопоставлена новая компактная реальность, как уже было показано в Какании Музиля. Моральные наблюдатели и моралисты обеспечивают друг друга уймой работы, стремясь нанести друг другу удар в спину: моральному наблюдателю напоминают о его наблюдательской морали, моралисту --- о его скрытой аморальности.

И это все? Допускают ли границы порядков только антитезу, только полемическую взаимодополняемость обращений к себе и к чуждому? Альтернативой, соответствующей второй форме обращения к чуждому, было бы возвращение к \textit{обращению к себе в обращении к чуждому}, с новой формой ответственности, позволяющей сочетать неизбежность требований с созданием наших собственных ответов. Как станет видно далее, эта не-возможность может также вернуть к требованиям чуждого, противоречащим нашим собственным идеям и желаниям и нарушающим общие правила. Сопротивление чуждого проистекает из отсутствия какого-либо эквивалента --- включая моральный,--- ускользающего из порядка. Для всякого порядка чуждое остается вторгающимся.

\chapter{Между пафосом и ответом}

Если предположить, что чуждое нарушает границы любого порядка, то на что похож опыт такого нарушения? Не стоит ждать, что сочетание здравого смысла и правил, намеренных действий субъекта, руководящегося правилами, наряду с консенсуальным согллашением между субъектами, может противостоять вызовам чуждого. Рассматриваемая здесь альтернатива предстает в форме феноменологии, основанной на пафосе и направленной на ответность, однако альтернатива эта уводит нас на задворки особого рода феноменологии и герменевтики, которые нацелены исключительно на интерпретацию смысла.

\section{В сфере смысла}

Подобно Фрейду, полагающему бессознательное основой психоанализа, можно положить интенциональность основой феноменологии. В точном определении интенциональность означает, что нечто проявляет себя как \textit{нечто},\footnote{Это ''как'' знакомо нам из других языков как \textit{qua, als, comme,} или \textit{come.}} что нечто подразумевается, дается, понимается или трактуется определенным образом: зеленая трава, пятно крови, вкус клубники, стол, Пегас, арабское число, любовное письмо, компьютерная программа, инфекция, нападение, террористический акт и т.д.. Формула \textit{нечто как нечто} означает, что нечто --- действительное, возможное или невозможное --- связано с чем-то другим (смысл, значение) и в то же время отделено от него. Реальность и чувство нельзя противопоставлять как характеристики или ценности. Подобно фуге, соединяющей несоединимое, миниатюра ''как'' обозначает разорванную связь. Мерло-Понти говорит о нулевой точке, стыке или складке, характеризующей возникновение смысла и Гештальта. ''Как'' не образует третьей сущности, перемещающейся между двумя исходными реальностями --- реальной и идеальной или физической и психической; скорее оно значит динамическую структуру, без которой буквально ничего проявляющего себя не было бы --- не было бы, значит, никого, кому оно является. Ничто не дается не будучи дано \textit{как нечто}, и никто не реагирует на него не реагируя \textit{как некто}. Опираясь на порождающий смысл аспект этого дифференцирующего процесса, я часто говорил об \textit{обозначающем различии}. По сравнению с основным характером этого дифференцирующего процесса, обращение к актам наделения смыслом предстает как специфическая интерпретация. Тот факт, что я замечаю оговорку, записку на двери или странный звук двигателя, еще не имеет характера приписываемого мной себе действия.

Если придерживаться подлинного дифференцирующего характера учения о значении, то обнаруживаются не только взаимосвязи между разными версиями феноменологии, но и связи с герменевтикой присутствия, интерпретацией традиций и текстов, семиотическими подходами, и, наконец --- что не менее важно --- с аналитической философией. Это делает идею смены парадигмы в феноменологии от философии сознания к семантике, предложенную Эрнстом Тугендхатом, гораздо менее правдоподобной. Несмотря на все разночтения между упомянутыми здесь альтернативами, на первый план выходит нечто общее, приводящее к философии смысла --- которая обычно заводит в тупик обычную эпистемологическую борьбу вроде таковой между реализмом и идеализмом, одновременно прекращая перетягивание каната между субъектом и объектом. Разделение на внешний и внутренний миры, введенное еще Джоном Локком и неоднократно после него повторявшееся, требует дополнительного третьего мира идей, поскольку оказывается конструкцией, покидающей почву опыта прежде ее обнаружения. Пронизанный интенциональностью, опыт не происходит ни внутри, ни снаружи. Так же и разделение эмпирических данных и общих идей оказывается подорванным. Как и в случае с формированием смысла и Гештальта, с самого начала опыт стремится к обобщению, не опираясь на данные в качестве основы. Присущие интенциональности ''как'' и ''каким образом'' подразумевают возможность повторения, и, следовательно, идеальность в генетическом и оперативном смысле, предшествующую всякой эйдетической или категориальной интуиции. Иначе говоря, интенциональный и дифференцированный опыт порождает множество горизонтальных и вертикальных опосредований, не полагаясь на заранее готовый разум или направляющего субъекта. Разум и субъект сами претерпевают генезис. Феноменология с концепцией горизонта смысла и опыта устанавливает совершенную пограничную концепцию и по-своему рассматривает вопрос об описании и преодолении границ без их устранения.

Обсуждаемое в первой главе мышление внутри и в пределах опыта также вытекает из базовой позиции, обозначенной как принцип феноменологии --- требуется лишь небольшой шаг, чтобы наделить \textit{нечто как нечто} коэффициентом контингентности. То что нечто предстает как нечто значит, что оно предстает \textit{только так и никак иначе}. С точки зрения теории Гештальта, каждый смысл, раскрывающийся как контекст ссылок --- привилегированный смысл: растение можно использовать как лекарственное или оставить как сорняк, нож можно использовать для еды или как оружие, денежный взнос может рассматриваться как пожертвование или как взятка, к иностранцу могут относиться как к просителю убежища или как к нелегалу. Как подчеркивал Ницше, привязанность к смыслу --- необратимая перспективность, но она же и потенциальная возможность конфликта. Есть смысл, но не какой-то отдельный смысл --- смысл развивается на фоне бессмыслицы, как показывает Мерло-Понти в своей \textit{Феноменологии восприятия} (Merleau-Ponty, 1945, 342; 2002, 345). Потому феноменология --- как и герменевтика --- повсюду приобретает индексальный характер. Происхождение системы Здесь-Сейчас, которую Карл Бюлер --- вдохновленный \textit{Логическими исследованиями} Гуссерля --- избрал в качестве центральной точки своей теории языка, не представляет собой комплекса принципов --- оно формирует поле, поле указаний, лежащее в основании символического поля с его расширениями из контекста и деконтекстуализациями. Помимо повседневных смысловых сочетаний, эта индексальность таит определенные исторические черты, свидетельствующие об эпохальных изменениях. Даже те события, которые мы обычно называем ''установлениями'' [\textit{Stiftungen}] --- рождение трагедии, демократии или геометрии --- имеют свои особенности, свои фазы инкубации и постановления. Их развитие не может быть подчинено телеологии, хотя близость понятия смысла к традиционному понятию цели или предназначения к этому манит.

Итак, я не вижу причин отказаться от этой плодотворной перспективы и занять другую. Поступать так было бы наивно: это подразумевает традиционное столкновение принципов или методов. Ориентация на смысл лежит в основе различных теорий значения, понимания и коммуникации и даже играет определенную роль в системно-теоретических подходах. Ее нельзя отрицать, не замалчивая. Но то, что нельзя отрицать, еще не окончательно оформлено и не гарантирует своей невосприимчивости к подрывным действиям, и что справедливо для смысловых формаций, то тем более справедливо для тех притязаний на достоверность, которые основаны на дискурсе и нормативно служат для фильтрации смысловых формаций. Речь никогда не о резких противопоставлениях или простых дополнениях: нас интересуют изменения, проявляющие свой устойчивый эффект подобно нюансам цветной живописи.


\section{Бездны претерпеваемого}

Сомнение в интенциональности должно вытекать из нее самой: как и в случае понимания или коммуникации, мы обречены на неудачу, хватки которой пытаемся избежать. Прямая конфронтация с ней означала бы возврат к лишенным смысла и норм данностям --- притворство, будто есть что-то незатронутое предписаниями смысла и норм. Предположение о таком ''чем-то'' рано или поздно попало бы в сети здравого смысла и было связано цепями аргументации. Следуя основной формуле, выбранной для этого анализа, мы должны понять --- как только внимательней рассмотрим возникновение смысла --- что есть ошибочные пути особого рода. То, что нечто проявляется как \textit{нечто}, не значит что это \textit{нечто}. Оно \textit{становится нечтом} через обретение смысла, и так становится высказываемым, доступным, повторяемым. Недостаточно предполагать происхождение смысла --- как если бы что-то просто выходило на свет из глубин опыта. Речь о достижении, порождении смысла, как говорят о генерации энергии. В этом гетерогенезе, проливающем также новый свет на гетерономию, вступает в игру нечто, что пока смысла не имеет, и вопрос в том, как понять этот \textit{гетеронт} как таковой.

Для прояснения позвольте привести отрывок из Роберта Музиля --- автора, не только близкого к Эрнсту Маху и вдохновлявшемуся Ницше, но и знакомого с ранними работами Гуссерля. В самом начале \textit{Человека без свойств} венская пара становится свидетелем события, которое мы обычно описываем \textit{как дорожно-транспортное происшествие} (так же как описываем некоторые метеорологические процессы и соответствующее им барометрическое давление, влажность и расположение солнца как \textit{великолепный августовский день}). Но вначале все не так. Все начинается с ''толпы'', с блокированного движения. Я слышу, как что-то тяжелое ''поворачивается в сторону, затем раздается звук резкого торможения, и то, что, по-видимому, является тяжелым грузовиком, застряло одним колесом на краю тротуара''. Рядом с ним стоит ''серый, как асфальт'' водитель и пытается объяснить, что произошло. Взгляды прохожих устремляются в ''глубину замкнутого пространства'', где жертва несчастного случая лежит как мертвая. Кто-то вызвал ''скорую'' --- ''экспертную, авторизованную помощь'', для которой экстренная помощь является нормальным явлением. Тем временем ''всезнайка'' объясняет какой-то женщине, что значит ''слишком большой тормозной путь''. Женщина благодарна за эти объяснения: ''для нее было достаточно того, что так это ужасное происшествие можно вписать в какую-то схему, превратить в техническую проблему, больше не касающуюся ее напрямую''. Это становится примером для ''социальных институтов'', которые так замечательно точны в своих услугах. ''У людей сложилось почти оправданное впечатление, что произошелшее было событием в рамках закона и правопорядка''. Комментарий разговорчивого мужчины и реакция его спутницы подтверждают это впечатление: '''Согласно американской статистике',-- отмечает мужчина,-- 'ежегодно на дорогах гибнет более ста девяноста тысяч человек и четыреста пятьдесят тысяч получают ранения'. 'Думаете, он мертв?' -- спрашивает его спутница, все еще питая неоправданное впечатление, что пережитое ею было чем-то исключительным'' (Musil 1978, 5-6).

Повседневное/неповседневное происшествие, которое здесь заканчивается, показывает, как трагическое событие превращается в статистическую норму, в ''упорядоченное'' событие: наделенное смыслом и соответствующее правилам, хотя все же оставляющее после себя субъективные ощущения в качестве прибавочного аффекта. Перед лицом растущей анонимизации кажущегося наиболее родным, это чувство напоминает личную реликвию: в конечном счете мы получаем ''переживания без переживающих'', предсказывает Музиль (1978, 150). Конечно, эту открывающую путешествие по Какании прелюдию можно было бы продолжить с помощью Фуко или Лумана, позволив, например, многофункциональной системе, называемой обществом, проявить себя со всеми ее органами чувств --- мы обнаружили бы различные коды и руководящие различия, будь то в медицинских тестах, юридических доказательствах или оценках действий, страховых агенствах, освещении инцидентов в средствах массовой информации --- когда уделяемое инциденту внимание зависит от общественной значимости жертвы и преступника,--- наконец, в литературном изложении, приобретающем эстетическую ценность в соответствии с разделением труда в этом производстве смысла, и все это не говоря уже о наших собственных феноменологических достижениях. Электрическая цепь смысла замыкается.

Безусловно, Гуссерль в \textit{Кризисе} высмеивает тех, кто приравнивает профессию феноменолога с таковой сапожника из пословицы, который всегда занят только своим делом\footnote{Let the cobbler stick to his last.} (Hua VI:140). Здесь можно было бы вновь процитировать Музиля, голосом главного героя смеющегося над ситуацией, когда ''птицеферма была создана ради возвышенных мыслей, называемых философией, теологией или литературой'' (Musil 1978, 358). Однако разделение профессий на специализации не должно немедленно приноситься в жертву какой-либо благородной цели: ни Гуссерль, ни Музиль этого не делают, еще меньше это делает Макс Вебер. Вполне возможно, что философ --- а следовательно, и феноменолог --- также сталкивается с определенными профессиональными требованиями, но он должен иметь способность и желание отразить профессионализацию образцовым образом. Музиль как литературный автор справляется с этой задачей: он не просто рассказывает о происшествии и вызывает тем сочувствие у читателя --- он ставит его перед фактом трансформации смысла значимого события. Есть нечто, что будучи хорошо подогнано и подготовлено по-другому, попадает в другие ''области смысла''. И все же, как возникает это ''незнамо-что'', что позволяет так легко представить себя в том или ином виде? Оно возникает как неожиданное событие, подобно corpus delicti ''ломает ряды'' и ''пересекает'' обычный ход событий, привлекая внимание, привлекая любопытных, разжигая желание увидеть, превращая наблюдателей в вуайерестов, возникает как то, что произошло с жертвой из-за ее ''беспечности''.

Несчастный случай, о котором говорится в романе --- прототип всего, что приходит нам на ум, что привлекает наше внимание, отталкивает, провоцирует, причиняет боль, заставляет думать и, в худшем случае, уничтожает нас как ''мыслящий тростник''. Здесь мы могли бы обратиться к феноменологии внимания, которая будет рассмотрена отдельно в одной из последующих глав. Что такое внимание? Обладает ли оно смыслом или это просто избирательная форма намерения? Так было бы только если бы внимание было сужено волюнтаристским образом или приравнивалось к лучу профектора, высвечивающего то, что всегда уже ожидает нашего взгляда в темноте. Однако привлечение внимания или осознания --- не те действия, которые можно совершать или от которых можно отказаться. Что-то может привлечь наше внимание как в виде легких, кажущихся безобидными, отвлекающих факторов, так и в форме сильных потрясений и опасностей --- как в описанном случае жертвы дорожного происшествия, которая следует своим намерениям, когда ее жестоко отрывают от них. Для обозначения этих --- не находящихся в нашем распоряжении, не ждущих разрешения или приказа, а просто происходящих с нами, одолевающих, волнующих, удивляющих нас, нападающих на нас --- событий, лучше подходит старый термин ''пафос'' (pathos), который, согласно высказыванию πάθει μάθος, говорит о познании через страдание, но не о познании страдания.\footnote{Старую доктрину катарсиса также необходимо бы снова подвергнуть сомнению --- вне зависимости от того, понимается ли катарсис как очищение аффектов или от аффектов. Процесс, понимаемый наполовину с медицинской, наполовину с моральной точки зрения, и ранее называвшийся ''очищением'', а сегодня --- ''проработкой'', слишком легко превращается в эмоционально управляемое упорядочивание претерпеваемого.} Такая перспектива освободит нас от нагромождения субъективных, личных эмоциональных состояний, проистекающих из ''дополнительной абстракции'' (\textit{Hua} VI:321; \textit{Crisis}, 226) и служащих для компенсации --- в форме психических данных --- утраченного при сведении мира к данным физическим. Но две вместе взятые абстракции не приводят к конкретизации: не хватает среды, в которой они могли бы ''срастись''.

Пафос не значит, что есть нечто воздействующее на нас, не значит он и что нечто понимается и интерпретируется как нечто иное. И меньше, и больше: это понятие избегает альтернативы причинности и интенциональности во всех их традиционных формах.\footnote{Что интенциональность и причинность не разделяет пропасть, стало ясно мне довольно рано (см. \textit{Der Spielraum des Verhaltens}, 1980, глава 4). Но потребовалось некоторое время для достаточно точной дифференциации концепции эффекта.} Полагайся мы только на причинно-следственную связь, нам пришлось бы судить с точки зрения наблюдателя, подобно полицейскому, берущему образец крови, или врачу, ищущему микробы, борющемуся с риском заражения и определяющему причину смерти. Если же только на чистые процессы постижения, интерпретации и анализа --- это значило бы, что мы уже вышли на уровень осмысленности опыта, где его переломные моменты уже скрыты за интерпретацией смысла.

Пафос --- это когда на нас воздействует \textit{что-то}, и ''откуда'' не основано на предшествующем ''что'' и не может быть преобразовано в последующее ''для чего''. Пафос может быть эпизодическим, хроническим, ''вулканическим'' или ''нептуническим'', может быть разной степени интенсивности. Противоположность пафоса --- не противоречащее смыслу и не бессмысленность в обычном понимании (возникающая из-за обманутых ожиданий), а апатия, когда происходящее не имеет значения и все погружается в монотонность не-различия, как для смешного человека Достоевского, для которого даже дотянуться до пистолета кажется непосильным усилием. Подобно Лейбницу, Гуссерль использует для таких случаев метафору сна: как только наступает ''аффективное облегчение'' (\textit{Hua} XI: 168; \textit{ACPAS}, 216), переживание погружается в сон. Не делает ли метафора сна уход из жизни слишком безобидным? Даже Платон идет еще дальше, сравнивая полную апатию с отсутствием желаний в камне (\textit{Gorgias} 492e).

\section{Претерпевающий и отвечающий}

Пафос нарушает ряд привычных различий. Прежде всего это касается различия между субъектом и объектом, мужду объективным событием и субъективным актом. Возьмем классический пример удовольствия и боли. Эти состояния или события нельзя отнести к вещам в мире: боль без ощущающего ее мыслима только как результат патологического расщепления, и даже такого рода явления предполагают некоторую связь между претерпевающим и его телесным состоянием. Удовольствие и боль также и не субъективные действия: будь они таковы, испытавший их сам бы нес за них ответственность --- они были бы приписаны ему, включены в его горизонт понимания. Пафос --- событие, но особо рода: оно происходит с \textit{кем-то}. Что-то происходит со \textit{самостью}, к чему \textit{я} имею то или иное отношение. ''От чего'' пафоса превращается в ''к чему'' ответа в реакции словами или действиями, отвержении, приветствии или выражении. Способность ''говорить о том, от чего я страдаю'' --- которую поэт называет даром божьим --- значит высказывание особого рода даже в повседневной жизни: это высказывание и действие, начинающееся не с самого себя, а \textit{откуда-то}, откуда несет следы влияния чуждого. Самость --- без которой никто не был бы собой --- может возникнуть только из открытости к всегда ускользающему чуждому. Именно это я называю ответностью. Сущность, в современную эпоху называемая ''субъект'', предстает сначала как претерпевающая и отвечающая так, что я оказываюсь вовлечен не как инициатор, но как подверженный определенным переживаниям, суб-ект в необычном смысле этого слова, подразумеваемом Лаканом и Левинасом. Наш ответ, как и пафос, находится по ту сторону интенциональности. Ответность выходит за рамки всякой интенциональности, потому что отвечание на происходящее с нами не исчерпывается значением, пониманием или истинностью нашего ответа. Все это не ограничивается аффективной подоплекой нашего когнитивного и практического режимов поведения, а касается самой сути этих режимов: всякий раз, когда в быту, политике, науке или философии возникает что-то новое, с нами что-то происходит, что-то ''осеняет'' нас. Даже история живет, потому что воспоминания пробуждаются, а не просто хранятся и обрабатываются как данные. Престиж компьютеров и роботов растет по мере того, как восприятия и практики становятся невосприимчивыми, охлаждаются и замораживаются, ограничиваясь достижениями, на которые не влияют достигающие. Аф-фекты, всегда тая в себе некий дисфункциональный динамит, превращаются в своего рода топливо, которое можно  пополнять и расходовать по нашему желанию. Это позволяет создавать чувствующие машины вроде ЭМО-машины Дитера Дернера, которой мы можем приписать эмоциональные движения без эмоций, свободную адаптацию Музиля.\footnote{Следует, однако, отметить растущее значение эмоций как в психологии, так и в нейрофизиологии после их долгого нахождения в тени более доступных для компьютерного моделирования когниций. Хьюберт Дрейфус уже давно привносит феноменологические аспекты в эту дискуссию, уделяя больше внимания телесной ситуированности, привыканию и неявному знанию, чем традиционным ролевым моделям.} Конечно, было бы нелепо обвинять машины в том, что, по сути, относится к управлению ими.

\section{По эту сторону добра и зла}

Еще одно могущее вызывать сомнение различие --- между вопросами факта и вопросами права, между ''есть'' и ''должно'', а также между бытием и ценностью. Происходящее с нами и разворачивающееся в различных красках --- ужасающих, удивительных, привлекательных или тревожащих,--- ни регистрируемый нами факт, ни обязательная к следованию норма, ни вытекающая из наших оценок ценность. Всякое обоснование, принятие или оценка уже запоздали: мы имеем дело с чем-то, что \textit{nolens volens} предполагаем и всегда предполагали, когда занимали определенную позицию к нему, утверждая или отрицая, выделяя или исключая. В этом отношении описанный несчастный случай парадигмален: он удивляет всех причастных, хоть механизмы преодоления и защиты и вступают в силу довольно скоро, оставляя от него лишь статистический разброс. Морализаторы любят делать вид, будто между делаемым и случающимся есть пропасть. Это тенденция топить претерпевание в вопросах вины, она началась с Иова, которому пришлось отбиваться от обвинений друзей как от назойливых мух. Проблема не столько в том, правы ли спорщики: даже в лучшем случае они \textit{всего лишь} правы. Происходящее с нами --- в большем и меньшем масштабе, приятное, печальное и ужасное --- демонстрирует форму неизбежности, которую древние греки называли \textit{tyche}, и которая заслуживает антиморалистического выплеска в \textit{amor fati} Ницше.

Гуссерль уделял особое внимание допредикативным процессам, относящимся к опыту, потому что возникающий порядок вещей предшествует любому явному ''да'' или ''нет''. Так же есть и донормативные процессы, которые, конечно, не просто часть ''раскрытия слов'', раскрытия возможностей смысла. Они говорят с нами и затрагивают нас. Они обеспечивают основу для нормативных суждений, но не наоборот. Упускающие из виду то, что заявляет о себе здесь и сейчас, будут отставать не только от своих же возможностей, но и от требований чуждого. Лучший способ избежать моральных испытаний --- избегать ситуаций, в которых они могут возникнуть. Описанный несчастный случай в этом отношении достаточно безобиден: ни у кого из участников не возникает проблем из-за него. Но так не всегда, а в некотором смысле даже всегда не так: даже ''сторонний наблюдатель'' может быть привлечен к ответственности за неоказание помощи. Описанный здесь способ видения и мышления приобретает надлежащий вес только когда пассивность называемого Гуссерлем пассивным синтезом воспринимается не как лишь предшествующий или подлежащий уровень активных чувственных образований, а как радикальная форма ''изначальной пассивности'', проистекающая из аффекта и так постоянно привносящая ''чуждые эго'' [\textit{Ichfremdes}] элементы, которую следует приписывать переживанию затрагивающего меня события [\textit{Widerfahrnis}]. Таким образом, мы достигаем точки, когда происходят события, для которых нельзя назвать достаточных условий возможности. Это относится ко всем постановлениям [\textit{Stiftungen}], открывающим новую сферу смысла --- будь то научные открытия, художественные инновации, политические или религиозные реформы или сдвиги в философской мысли. То, что возникновение порядков ретроактивно проявляется в нашем инвентаре порядков, не значит, что инвентарь этот обеспечивает достаточный запас: драматические события происходят не только в истории, но и в повседневной жизни, где удивляют нас как \textit{objets trouvés} и \textit{personnes trouvées} и ведут свою скрытую игру в духе шутливости, \textit{jeu d’esprit}. Ключевые события любого рода заставляют противопоставлять принципу достаточного основания принцип \textit{не}достаточного основания: от случая к случаю все попытки рассуждения и обоснования наталкиваются на непреодолимые пределы. Как не поддающееся компетентному описанию и объяснению можно, тем не менее, заметить и выразить? Даже молчанию нужно красноречие, чтобы что-то прояснить.

\section{Темпоральный \textit{диастаз}}

Всякий раз, когда нечто случается с нами, мы склонны --- от этой склонности трудно избавиться --- приписывать одну часть внешним причинам, а другую --- спонтанным актам свободы, как если бы интересовались соотношением между молотом и наковальней. Если что и удерживает нас от такого, так это подлинная темпоральность, определяющая двойное возникновение пафоса и ответа с начала и до конца.

Возьмем еще раз приведенный выше пример дорожно-транспортного происшествия. Пара прохожих на месте происшествия, делающих несколько ободряющих замечаний, появляются не сразу: ''Еще мгновение назад что-то было не так''. Значит ли это, что наблюдатели прибыли слишком поздно? Когда им следовало прибыть, чтобы поймать событие с поличным? Может, им следовало ждать его, как заката? Но тогда не произошло бы ничего необычного: пресловутый \textit{nil novi sub sole} должен был бы повлиять на само солнце, если бы привычка победила. Но как быть с преступником, пытающимся объясниться, и с жертвой, чья беспечность и привела к инциденту? Очевидно, событие произошло слишком рано, ответ последовал слишком поздно --- но слишком рано и слишком поздно в каком отношении? Определенно, их нельзя сопоставить с самим инцидентом --- он происходит именно в этом временном промежутке и нигде больше. Темпоральность, о которой речь, даже не принимается во внимание, если мы заранее вступаем в диалогическую плоскость, где свой и чуждый вклад синхронизируются посредством взаимности и обратимости перспектив: в диалоге никто и ничто не приходит слишком поздно, ведь всегда можно наверстать упущенное на основе общих предпосылок.

Временной сдвиг, возникающий от предшествования пафоса и отсрочки ответа и разделяющий гомогенный диа-лог на гетерогенный, я называю \textit{диастаз} --- изначальное расщепление, создающее контекст, пусть и разорванный. Предшествующий пафос и отложенный ответ должны мыслиться вместе, но только через непреодолимый разрыв, что требует творческого ответа. Происшествия не только ведут нас думать --- они заставляют нас думать. Разрыв, о котором речь --- определяющий для события в той же мере, в какой перспектива определяюща для восприятия: Гуссерль утверждает, что даже всевидящий Бог имел бы перспективу. Любые попытки устранить это напряжение --- путем освобождения пафоса от ответа или наоборот --- означают ступание на путь фундаментализма или конструктивизма, где одна сторона выигрывает от минимизации другой, противоположной. Опыт, начинающийся с событий, требует другого языка и другой логики. Здесь мы сталкиваемся с любопытным случаем: эффект \textit{sub specie patientis} предшествует его причине. Линейная временная схема, на которой классическая физика основывает свою концепцию причинности, терпит неудачу, если объект воздействия задним числом создает историю, возвращаясь к прошлому (см. \textit{Hua} X:54). Пафос и отклик не следуют друг за другом подобно двум событиям --- это даже не два разных события, а одно переживание, сдвинутое по отношению к самому себе: подлинный временной лаг. Эту гетерохронию не устранить синхронизацией: синхронизация уже предполагает темпоральную \textit{distentio}, расширяющую любое значимое \textit{intentio}. Дух пафоса, говоря словами Паскаля --- ''хромающий дух'', или, умеренней, ''запаздывающий дух''. Поэтому мы изначально сталкиваемся с пафосом не как с чем-то, что подразумеваем, понимаем, осуждаем, отвергаем или утверждаем --- он формирует время и место, \textit{из которого} мы будем делать это всё, отвечая на него. Все, что со мной происходит и на что я отвечаю, не имеет смысла как такового и не подчиняется какому-либо правилу. Настаивание на истинном удивлении, оправданном изумлении или неуместной тревоге накладывает на переживание критерии, вступающие в силу только в его обработке. Удивление остается удивлением, а тревога --- тревогой, если только мы не подвержены невежеству, наивности или фантазиям, которые растворяются по мере овладения опытом.

Тем не менее, нарушающее порядок интегрируется в него, будучи названо, классифицированно, датированно, локализованно и объяснено. Это возвращение, относящееся к эффектам событий и обеспечивающее им определенную продолжительность, выходит на первый план в описании происшествия Музилем. Он делает даже большее: приостанавливает эту нормализацию, делая ее последствия заметными, придавая ей иронический оттенок и так лишая ее всякой видимости престижа. Читая его описание, мы ощущаем, как событие скрывается за ''покровом смысла'', постепенно обретает осмысленность, даже чувственный интерес, что дает почву для спекуляций и распространения. Сочувственный тон иронии здесь менее важен, чем ее косвенность, указывающая на несовпадение между сказанным и понятым о событии и тем, что произошло на самом деле. Пафос, не исчерпываясь определением \textit{pathetique} и даже не завися от нее, сталкивает нас с избытком, который никогда не будет исчерпан. Это нечто, что следует обозначить как бессмысленное и бесцельное --- оно разрывает сети смысла, нарушает систему правил и так деконтекстуализирует событие. Оно не-посредственно только потому что прорывается через посредничество. Для тематизации этих проявлений патетики потребуется ответностное \textit{epoché} или ответностная редукция, отслеживающая о-смысленные, правило-ведомые и ценностно-релевантные высказывания вплоть до вызывающего их ответ (см. \textit{Antwortregister, 195}). Высказывание и деяние исчезают в осмысленности и упорядоченности сказанного и сделанного, исчезновение это можно остановить только непрерывным перевысказыванием (\textit{redire}) и высказыванием-против (\textit{dédire}), как того требует Левинас (1974, 8-9, 197-98). Логичность смысла и канон правил так не игнорируются, а нарушаются очевидностью происходящего с нами и заставляющего нас отвечать. Пафос, традиционно ассоциируясь с нелогичным и иррациональным, присутствует в самом логосе. Генеалогия логики или морали, стремясь быть чем-то большим, чем история идей или нравственности, возможна только если оставить прочные основы здравого смысла позади, вернув разуму и свободе их исходную природу --- природу бездны. Для феноменологии это значит необходимость обратиться против самой себя, противостоять эйфории смысла, притупляющей ее так же как и прочие философии смысла. Здесь я вижу основные проблемы феноменологии, и они не касаются устаревшего наследия философии сознания, могущей мыслить бытие-сознание только в фреймворке независимой области сознания.

\section{Проверка практикой}

Можно ожидать, что феноменология, основанная на пафосе и ответе, расставит свои акценты в текущих дебатах. Ниже мы подробно рассмотрим, как это можно представить в отношении некоторых соответствующих тестовых примеров.

Насилие в личной и коллективной истории рассматривалось иудейско-христианской традицией как зло, и в последнее время от его рассмотрения слишком легко переходят к общему поиску смысла или к морально-юридической оценке. Насилие становится заметней, когда с самого начала воспринимается как пафос, как травма, которая случается с кем-то или причиняется кому-то. Насилие тогда предстает как чуждая материя, разрушающая существующие смысловые структуры.\footnote{Мои размышления об ''апориях насилия'' см. в Dabag, Kapust, Waldenfels 2000.} Это не значит, что здесь не за что ухватиться, понять, осудить, исправить или предотвратить, но суть всех этих попыток недоступна псевдорационализации и одностороннему морализаторству. Даже компенсация и возмещение ущерба не залечивают ''шрамы истории''. Травмы принимают крайние формы при блокировании или сведении на нет всех попыток ответа. За всяким отказом от дискурса стоит молчание, погружение в тишину, и мы принижаем значение этого переживания, возлагая ответственность за подобные нарушения порядка на медицину. В сфере травматизации своего рода отсрочка, après coup, характерная для всего патологического, проявляется в патологической форме фиксации на происходящем с последующей блокировкой ответа. На этом этапе феноменология времени соединяется с психоаналитической археологией до-времени. Наконец, даже историография не остается в стороне от этой проблемы: события прошлого забываются, подавляются, и их необходимо пробудить, прежде чем можно будет приступить к архивной, монументалистской или нравственной работе по запоминанию, что приведет к формированию так называемой культуры запоминания. Ницше убедительно показал, насколько тесно связаны боль и память, перенесенное и запомненное. В этом отношении настойчивое стремление Якоба Буркхардта принять страдающего человека в дополнение к человеку стремящемуся и действующему в историологии сейчас стало уже почти общественным императивом, тогда как в философии по разным причинам по-прежнему преобладает активизм односторонних теорий действия.

Биоэтические дебаты могут принять иной оборот, если в случае рождения и смерти, а также в случае неизлечимой и опасной для жизни болезни мы начнем с воплощенной самости, имеющей свою пре-, пост- и глубинную историю.\footnote{Более подробное рассмотрение см. в \textit{Bruchlinien der Erfahrung}, главы 8, 11-13.} Эта самость с самого начала была бы больше, чем лишь кому-то приписываемые определенные роли, права и способности. Чудовищные идеи, подобные идее о кучке клеток, неделенных достоинством, были бы тогда отброшены как пример той самой двойной абстракции, о которой говорилось выше. Происходящее с самостью на уровне тела множится в события событий, когда мы рассматриваем вовлеченность в окружающий мир, т.е. вовлеченность всех тех, кто, рожая, ожидает неожиданного, или, в случае смерти, скорбит и прощается с ускользающим от воспоминаний. Без таких пре- и постисторий и сопутствующих им прерываний, жизнь была бы продуктом наших действий, даже когда \textit{poiema} усиливалась бы \textit{autopoietic} образом. Старое определение жизни как самодвижения не вполне ложно, но недостаточно для мышления о само-предшествии, которое включает патологическую форму \textit{гетерокинеза}.

Дискуссия о чуждом получит шанс освободиться от метаний присвоения-отбрасывания, от попыток интегрировать чуждое и уступить ему, только если чуждое будет рассматриваться с точки зрения пафоса --- как нарушение или помеха, как при воздействии от чего-то не закрепленного за каким-либо объектом или значением. Это верно как для межличностных отношений, так и для внутри- и межкультурных обменов. Переживание чуждого не начинается ни с хороших, ни с плохих намерений: оно разрушает ожидания смысла и порядка, на которых эти намерения основаны. Пафос --- не просто ненамеренное, но не могущее быть намеренным. С философской точки зрения, чуждое --- нечто, которое, несмотря на все потенциальные возможности личностного, историко-культурного или даже трансцендентального характера, оказывается не-возможным, нарушением существующих возможностей или вызовом им. Это направление мы находим у таких авторов, как Бергсон и Джеймс, где оно объявляет себя ''радикальным эмпиризмом'', несколько иначе оно реализуется Левинасом и Делезом. Сколь актуален бы ни был вопрос о возможностях обусловливания опыта, он не является надежным якорем против его перипетий.

\chapter{Ответ чуждому}

Вопрос об опыте чуждого неразрывно связан с вопросом о том, как мы встречаем чуждое. Вопрос этот отражается в мотиве ответности, который немыслим без этического измерения. Однако руководимая этим мотивом ответностная этика идет глубже, чем философия морали, с самого начала опирающаяся на заповеди, права или ценности. Зов чуждого, в отличие от не касающегося нас факта, не безразличен, однако, это не значит его подчинение определенным стандартам достоверности. Скорее, без элемента аморальности всякая мораль превращается в мораль стада.

\section{Ответность}

Феномен чуждого, которое проявляет себя, ускользая, можно охарактеризовать как \textit{гиперфеномен}. Соответственно, Гуссерль характеризует чуждое как ''поддающуюся проверке доступность изначально недоступного'' (\textit{Hua} I:144), и в одном из своих томов, посвященных феноменологии интерсубъективности (\textit{Hua} XV:631), он столь же парадоксальным образом говорит о ''доступности в подлинной недоступности, в режиме непостижимости''.\footnote{''Zugänglichkeit in der eigentlichen Unzugänglichkeit, im Modus der Unverständlichkeit''.} На социальном уровне мы сталкиваемся с соответствующей структурой ''принадлежности в не-принадлежности'': каждый, кто принадлежит к семье, народу, касте, религиозному сообществу или культуре, никогда не принадлежит к ним полностью. Отстраненность, дистанцированность, отдаленность, а также моменты одиночества и бытия-вне-места --- на которые феноменологи часто ссылаются в анализе опыта чуждого --- не означают ослабления этого опыта, напротив, они относятся к его сути. В основе любой феноменологии чуждого лежит понимание, что переживание чуждого не значит недостаток, так же как недостатка нет и в переживании лежащего в прошлом или будущем. Радикальный характер чуждого не означает, что это нечто совершенно отличное от своего и привычного --- он означает, что его нельзя ни вывести из своего, ни отнести к общему. Как нечто самоустраняющееся, чуждое не только гиперфеномен, но и изначальный феномен, подобный феномену контраста (см. \textit{Hua} XI:138). Именно поэтому каждый опыт чуждого находится по ту сторону смысла и правил, то есть сравнивается со смыслом, \textit{в соответствии с} которым мы понимаем нечто и самих себя, и сравнивается с правилами, \textit{которыми} мы руководствуемся, так или иначе относясь к кому-то или чему-то. Общие характеристики \textit{интенциональности} и \textit{регулярности}, от которых зависит возникновение общего мира, не заменяются \textit{ответностью}, но, безусловно, превосходятся ей. Ответность предшествует ответственности за все что мы делаем и говорим.\footnote{Термины ''ответность'' (responsivity) и ''ответ'' я позаимствовал из области медицины (а именно, у Карту Гольдштейна) и бихевиористской психологии поведения (см. \textit{In den Netzen der Lebenswelt}, 1985, 132-33, \textit{Antwortregister}, 1994, 457-61). Эквивалент можно найти у Михаила Бахтина, который в своей ''теории романа'' использует необычное выражение \textit{ответность} (Бахтин 1979, 233), на английском буквально ''answerability''. Наконец, во французском языке есть особая близость между ответом и принятием ответственности в выражении \textit{respondeur de ...}, что значит, что человек берет на себя ответственность за что-то или кого-то. Левинас, а также Бланшо, преследуют эту связь. Я ссылаюсь на главу 19 \textit{Deutsch-Französische Gedankengänge} под названием ''Antwort der Verantwortung'', которая посвящена Левинасу и берет свой девиз от Бланшо: ''Répondre de ce qui échappe à la responsabilité''.}

Здесь мы видим следы другой --- не ограничивающейся сферой смысла --- феноменологии, ее логос демонстрирует черты изначальной гетерологии. Нарушение сферы интенционального или управляемого правилами смысла происходит в ответ на \textit{требование чуждого --- требование, не имеющее смысла и не следующее правилу}, но разрушающее привычные смысловые образования и правила, провоцируя тем самым создание новых. \textit{Что} я отвечаю обязано своим значением вызову, \textit{на который} я отвечаю. Чуждое, предстающее перед нами как зов чуждого или взгляд со стороны чуждого, теряет свою чуждость, если \textit{ответностное различие} между тем, чему мы отвечаем, и тем, на что мы отвечаем, заменяется интенциональным или управляемым правилами смысловым процессом. Ответностное различие исчезает за \textit{значностным} или \textit{герменевтическим различием}, в котором нечто вопринимается или понимается как нечто, исчезает оно и за \textit{регулятивным различием}, в котором нечто рассматривается в соответствии с нормой. Феноменологическое, герменевтическое или регулятивное ''как'' скрывает опыт чуждого, когда понимание и коммуникация отрицают их ответностный характер.

Чуждое как таковое требует адекватной формы феноменологии, начинающейся с бросающего нам вызов, взывающего к нам или ставящего под сомнение наши возможности отчуждающим, шокирующим или поражающим образом, прежде чем мы вступим в нашу собственную хотим-знать и хотим-понимать ситуацию. Пафос чуждого превосходит его сомнительность. Однако это не тот случай, когда мы заменяем традиционное первенство сомнения, возникающее как либо непосредственно соответствующий стремлению к знанию вопрос о факте, либо косвенно выводящий нас на тот же путь интерсубъективный вопрос. Простые развороты, как правило, не ведут далеко, потому что остаются привязанными к тому, что меняют. Что нам нужно, так это смещение акцента и новая --- открывающая новые пути --- ориентация. Мы довольно долго пытались показать, что направленный на реализацию намерений и правил и симметричное уравнивающее распределение ролей вопросов и ответов, круг вопросов и ответов --- не самоочевиден. Очевидно тогда, что классический диалог, с помощью которого любая традиционная философия поднимается и падает, уже предполагает асимметрию призыва и ответа, а также разрыв между бытием-призванным и отвечанием.

Все выходящие за рамки здесь сказанного речи содержат обещание, невыполнимое путем консенсуса в очередности хода и согласованности регулируемого поведения. Часто цитируемый Гельдерлином ''разговор о том, кем мы являемся'' возникает из отстраненности чуждого, чей призыв предшествует любому партнерству.

\section{Зов и ответ}

Зов Другого разрывает целенаправленный круг интенциональности и регулятивный круг общения и проявляет в себе чуждое \textit{in actu}. Такого рода призыв --- или, говоря по-немецки, \textit{Anspruch} --- означает сразу две вещи: адресованное \textit{кому-то} обращение и претензию \textit{на что-то}. Особенность зова Другого --- переплетенность в нем обеих этих форм обращения, призыва и требования. В полученном мной призыве от меня что-то требуется. Ситуативно воплощенный призыв предшествует любому моральному или юридическому требованию: вопрос о законности требования предполагает уже воспринятость его как вызова. Здесь мы подходим к области, расположенной \textit{по эту\footnote{Ориг. hither.} сторону добра и зла, правильного и неправильного}.

Мораль здесь показывает нам свое слепое пятно: все попытки обосновать мораль предполагают фактические утверждения, большие чем просто факты. Случай, когда у меня спрашивают дорогу или мое имя, становится фактом только тогда, когда к нему относятся как к факту, могущему произойти в результате наблюдения, утверждения или пересказа, как у меня что-то спрашивали. Но то, что \textit{становится} фактом, \textit{не является} им. Фактуализация чуждого вопроса не может помешать мне быть призванным этим вопросом. Кроме того, отвечание начинается не с разговора, вовсе нет --- с разглядывания и выслушивания, что в некоторой степени неизбежно. Повелительное ''Слушай!'' не услышать, не вслушиваясь в него. ''Не слушай меня!'' приводит к знаменитой двойной привязке: никакой ответ не будет верным. Даже попытка не слушать [\textit{Weghören}] предполагает определенное слушание, а попытка не смотреть [\textit{Wegsehen}] --- определенное смотрение, так же как и неуважение Ницше включает в себя некоторое уважение.

Ответ, соответствуя двоякому требованию, также принимает двоякую форму: требование соответствует ответу. С точки зрения теории речевых актов, это связано с подходящим \textit{содержанием ответа}, заполняющим пробел в пропозициональном содержании вопросов или запросов. Такой ответ восполняет нехватку, но само событие ответа этим никак не исчерпывается: обращение ко мне соответствует \textit{ответу}, который не заполняет пустоту, а \textit{выходит навстречу приглашениям и зовам Другого}. Отвечающий дает не то, что у него есть, а то, что изобретает, отвечая.

Отвечание или отказ от него происходят на уровне речевого \textit{события}, а не содержания. Отвечание здесь не исчерпывается ответом, как и отрицание отличается от отрицательных суждений, а стремление к знанию --- от невежества. Данный ответ может быть передан машине, которая соответствующим образом отреагирует на запрос. Это не считается предоставлением ответа, за исключением случая предварительно созданного и теперь воспроизводимого из данных, диалога. Однако машина, которая перестает работать, не отказывается служить. Возьмем, напротив, Бартлби из одноименной повести Мелвилла. Своим стереотипным повторением ''Я бы предпочел не делать этого'' он не только отказывается служить, но и доводит отказ до крайности, не отвечая и обрывая всякую коммуникацию. Он, как мы бы сказали, невосприимчив к любезности, расточаемой ему его работодателем, филантропически настроенным юристом. Отказ отвечать, уклончивая формула которого отодвигает на второй план молчание, указывает на бездны дискурса, необъятные для фактических разногласий. Потому отвечание --- нечто большее, чем смысл-направленное или правило-руководимое поведение. Дополнительный признак удвоения ответа в событие ответа и содержание ответа --- возможность ответить встречным вопросом, т.е. ответом, который не ответ в смысле предоставления запрашиваемой информации. \textit{Keine Antwort ist auch eine Antwort}, как мы говорим на немецком, или: отсутствие ответа --- тоже ответ.

Отвечание на зов чуждого не ограничивается устными высказываниями. Во ходе выполнения запроса речь чуждого и мои действия часто совпадают, например, когда я делаю то, о чем меня просят. Более того, можно спровоцировать ответный взгляд или ответить на взгляд. Встречные взгляды --- часть повседневной жизни в больших городах. Отвечание воплощает дух чувств, простирающийся от ритуалов приветствия до игр влюбленных. В конце концов, древнее определение ''Человек --- животное, которое обладает речью или разумом'' можно переформулировать в ''Человек --- животное, которое отвечает''. Необходимо пересмотреть разницу между человеком и животным и человеком и машиной.

\section{Аспекты логики ответа}

Ответность как основная черта человеческого поведения требует особой \textit{логики ответа}, существенно отличающейся от логики интенциональных актов, логики понимания или логики коммуникативного действия. Это приводит к соответствующей форме рациональности --- рациональности, возникающей из самой ответности.\footnote{Впервые я развил эту идею в обсуждении открытых форм контроля в ходе собеседований и интерактивных мероприятий (см. \textit{Ordnung im Zwielicht}, 1987, 47).} Далее я очерчу четыре аспекта этой своеобразной логики и в то же время покажу, как традиционные темы выглядят в свете ответности.

(1) Зов Другого, будучи нечтом большим, чем часть целого или частность общего, приобретает особую форму \textit{уникальности}. Уникальность эта действительно проявляется во множественном числе, но так, что не позволяет провести различие между частным и общим. Такая уникальность не значит, что нечто происходит только раз, как, например, звук или преступление, означает она и не просто нечто классифицируемое как одно среди многого. Здесь не находящаяся на нижней границе всеобъемлющей универсальности и потому невыразимая, индивидуальность. Здесь уникальность событий, проявляющаяся как таковая при нашем отклонении от привычных событий и открытии нового способа видения, мышления и действия. В жизни отдельных людей, а также народов и культур есть ключевые события [\textit{Schlüsselereignisse}], ''которые не забываются'' потому что вводят символические порядки, устанавливают смыслы, создают истории, провоцируют ответы и создают новые обязательства. Французская революция в этом отношении была --- по крайней мере для участвовавших в ней европейцев --- не просто одной из многих революций, и благодаря этой уникальности она становится кристаллизующим ядром для различных мифов и обрядов. Только в глазах стороннего наблюдателя Французская революция стоит наряду с Американской, Русской или Китайской революциями, подобно как для повзрослевшего ребенка его мать уже воспринимается наряду с другими женщинами, а место рождения --- наряду с другими местами.

(2) Аналогично призыв Другого отвергает разделение между фактами и нормами, ''есть'' и ''должно'', доминировавшее в области практической философии со времен Юма и Канта. Призыв, дающий возможность что-то увидеть, услышать, подумать, почувствовать, возникает с некой \textit{неизбежностью}, \textit{ne-cessitudo} в буквальном смысле этого слова, не вытекающей из универсальных законов, а, как практическая необходимость, относящейся к сущностным предпосылкам нашего социального существования в мире. Неизбежность эта означает, что как только зов Другого дошел до меня, я не могу на него не ответить, так же как, согласно Полу Вацлавику, не могу не общаться --- отсутствие ответа это тоже ответ. Двойное отрицание --- не могу не отвечать --- известно из определения необходимости из модальной логики. ''Необходимо, чтобы p'' может быть переведено с помощью оператора возможности M как ''не M (не-p)''. Двойное отрицание относится к требованиям, могущим быть понятыми только косвенно, а именно как необходимые условия для получения определенного опыта без возможности получить его позитивно. В аналогичном ключе Кант говорил о ''факте разума'', Гуссерль --- о ''необходимости факта'' или ''абсолютном факте'', Сартр приписывал cogito необходимость свершившегося (\textit{nécessité de fait}). По своей радикальности зов Другого напоминает такие традиционные примеры, как стремление к счастью, движение к самосохранению, категорический императив и человеческая свобода, которые --- если существуют --- не зависят от нашего выбора, не нуждаются в обосновании. Согласно Платону, мы не руководим своим стремлением к счастью, так же как, согласно Канту, от нас не зависит, слышим мы голос морального закона или нет.

То, с чего начинаются --- и всегда начинались --- наш дискурс и речь, нельзя наблюдать, оценивать или направлять подобно нечту у нас перед глазами или в руках. Оно проявляется только когда мы что-то говорим или делаем. Оно зависит от косвенной манеры общения, которая связана --- и остается связанной --- с молчанием. Это верно и для влечения, исходящего от чуждого вообще и чуждой культуры в особенности. Конечно, как только достигнут этап ''культурного диалога'', между культурами происходят процессы взаимного изучения. Влечение это столь же мало контролируемо, сколь удивление или любовная мания. Волнениям такого рода можно либо поддаться, либо заранее избегать их: от них нельзя избавиться, они становятся как воздух.

(3) В предыдущей главе мы уже видели, что уникальные события не только приводят к неизбежному вызову, но и провоцируют недостижимую отсрочку [\textit{Nachträglichkeit}]. Логика ответа немыслима без такого рода темпоральности, ранее обозначенной как диастаз. Подлинная отсрочка ответа подрывает первичность некоего изначального присутствия. Присутствие --- не ничто, как полагают многие сторонники постмодернизма, хотя оно и не самодостаточно. Ответ происходит здесь и сейчас, но начинается он в другом месте. Деррида утверждает, что отсрочка эта приводит к появлению оригинала только ''путем замены'', те есть поддерживается дополнениями, тянущими за собой бесконечную цепочку ''оригинальных повторений''. Такую же отсрочку следует отнести и к тем травмирующим происшествиям, которые схватываемы только в их последствиях, как в рассказе Фрейда о детстве Человека-волка. Говорить здесь о простом коммуникативном искажении --- значит преуменьшать важность поставленного здесь на кон. Можно упустить из виду, что создание коммуникативного поля само по себе не происходит посредством коммуникативного соглашения --- так же как и Основной закон страны не принимается в соответствии с ним самим и соблюдение его не включено в разделы Конституции. Основание порядка --- событие, которое не часть этого порядка. В этом смысле каждое рождение, открывая новый мир, становится нечтом вроде перерождения, поскольку понятно только в ретроспективе. Свобода потому означает не способность начинать исключительно с себя --- это способность начинать где-то в другом месте. Тот, кто верит, что может начать с себя, лишь повторяет уже существующее, что значит не начинать. Ответить --- значит отказаться от первого, и, значит, от последнего, слова.

(4) С темпоральной отсрочкой требования и ответа рука об руку идет неустранимая асимметрия, выводящая из равновесия традиционный --- ориентированный на достижение общих целей и следование общим правилам --- диалог и оставляющая без внимания моральные требования равенства вроде тех, что содержатся в Золотом правиле. Как показывает Левинас, асимметрия эта зависит не от неравномерного распределения правил в текущем диалоге --- она зависит от несовпадения призыва и ответа. Между вопросом и ответом так же мало общего, как между просьбой и ее выполнением. Они сталкиваются подобно взглядам.

Принятие зова чуждого и дарование ответа превращаются во взаимное отдавание и принятие только когда свое и чуждое рассматриваются третьей стороной --- проводящей сравнения и в случае конфликта создающей баланс. Точка зрения третьей стороны, гарантируя закон и справедливость, в определенном смысле необходима. Поскольку в ходе обсуждения и действий повторяются формы и применяются правила или законы, в игре всегда присутствует какое-то третье лицо или инстанция. Однако, подчиняя призыв чуждого универсальному закону и так \textit{уравнивая неравное}, справедливость всегда содержит элемент несправедливости. Попытка создать определенную симметрию между своим и чуждым и так сделать их равными в конечном счете была бы подобна попытке уравновесить настоящее и прошлое, бодрствование и сон, жизнь и смерть, как если бы можно было по своему желанию переступить границу, отделяющую от другого. Чуждое этого не допускает: оно ведет себя как приходящие нам в голову идеи, навязчивые идеи, преследующие нас, как сны, от которых никогда полностью не просыпаемся. Оно приходит из необратимого ''однажды'' и необратимого ''где-то не здесь''.

Если отделять отвечание от заданного смысла и существующих правил, то нужно различать преимущественно \textit{повторяющееся} или \textit{репродуктивное} отвечание и преимущественно \textit{инновационное} или \textit{продуктивное} отвечание. Обычное, нормальное отвечание имеет смысл и следует определенным правилам. Однако это не верно для отвечания на непредвиденные, нарушающие существующий порядок и меняющие условия для понимания и коммуникации, вызовы. Но когда бы и где бы ни нарушался порядок вещей и слов, между провокацией чуждого, с которой мы сталкиваемся, и нашей собственной продукцией, возникает пропасть. Здесь парадокс творческого ответа, напоминающий парадокс творческого выражения Мерло-Понти.\footnote{См. ''Das Paradox des Ausdrucks'' в \textit{Deutsch- Französische Gedankengänge}.} Ответ творческий несмотря на то, что он --- ответ. Призыв не относится к порядку, объединяющему или подчиняющему ответ, он становится призывом \textit{в ответе}, который провоцирует и которому предшествует. Отвечание так проходит по небольшому гребню, отделяющему зависимость и уступчивость от произвола и своеволия. Тому, кто ждет готовых ответов, нечего сказать, ведь все \textit{уже сказано}. И тому, кто говорит, не отвечая --- тоже нечего \textit{сказать}, потому что ему нечего сказать. Мы изобретаем то, что отвечаем, но не то, на что отвечаем и что придает вес нашим словам и действиям.

\chapter{Телесный опыт между самостью и инаковостью}

Телесность и чуждость тесно связаны. Чуждость проявляется в плоти как отсутствие в плоти и кости (\textit{absence en chair et en os}), используя формулировку Сартра, намекающую на Гуссерлевское присутствие в плоти [\textit{leibhaftige Gegenwart}] воспринимаемого объекта. Телесное существо никогда не представлено себе полностью. Загадка чуждости, таким образом, усугубляется загадкой тела, и мы оказываемся снова втянутыми в картезианские и посткартезианские авантюры современности.

\section{Загадка тела}

Судьба современной эпохи глубоко определена совместностью возникновения и синергией математизации природы и господства эго. В одновременности этой все относящееся к нашему телесному существованию оказывается в тени дважды: из-за автономии субъекта и из-за измерения природы. И нет лучшего свидетельства этого процесса, чем мысль Декарта. Речь о мыслящем эго, которое относится к вещам: в некоторых вещах я узнаю себя как \textit{res cogitans}, другие проявляются как Другие, как и я, мыслящие, но большей частью это лишь физические примеры \textit{res extensa}. Но остается проблема оправдания такого вот особенного тела, обозначенного как мое (\textit{corpus meum}) и обнаруживающего другие, принадлежащие Другим, тела. Ирония истории заключается в том, что наше тело одновременно и скрыто и открыто: наши тела кажутся отделенными от нас самих, хотя в определенной степени принадлежат нам, особенно когда заставляют нас страдать.

Картезианский дуализм подрывается при первом же пересмотре, привносящем чувственный опыт и языковые выражения. Конечно, именно я отношусь к вещам, к самому себе и к другим, но делаю я это только благодаря своему постоянно вмешивающемуся телу. Никакие манипуляции невозможны без рук, никакое общение невозможно без рта и ушей, нет ощущений без кровяного давления и сердцебиения и т.д.. Но этот процесс конкретизации останавливается на полпути: все еще предполагается кто-то или что-то, что включено и воплощено, не являясь вполне телом. Значит, выход есть: благодаря всепроникающему Духу, разве нельзя мне возвращаться к самому себе и в то же время достигать тотальности существ, интернализуя изначально принадлежащее внешнему тела? Я так начинаю со своего ограниченного ''я'' и заканчиваю всеобъемлющим целым. Нелегко сказать, кто дальше от истины: Декарт, придерживающийся дуалистической позиции и настаивающий на разрыве между разумом и телом, или Гегель, несмотря на всю рассеянность, придерживающийся монистической --- интегрирующей тело в целостность Духа --- позиции. Отголоски этих старых дискуссий звучат сегодня даже в науках о человеке и жизни. Даже нейрофизиология настаивает на различии между взглядами от первого и третьего лица, но --- несмотря на полезность этого различия на методологическом и экспериментальном уровнях --- остается открытым вопрос, как определить зловещий X, проявляющийся, как предполагается, с двух разных точек зрения. С другой стороны, любой вид монизма --- будь то физический или биологический --- сам по себе вызывает проблемы, когда сталкивается с существами, которые не только находятся под нашим наблюдением и нами обсуждаются, но и сами говорят с нами и смотрят на нас.

Есть, впрочем, другой вид пересмотра, выходящий за рамки этой картезианской или полукартезианской проблемы: не верно ли скорее, что в определенном смысле я \textit{и есть} мое тело и ты \textit{и есть} твое тело, как утверждают некоторые феноменологи? Действительно ли я \textit{имею} боль и воприятия как если бы был собственником своего опыта? Если мы предположим, как это делает Гельмут Плеснер, что бытие-телом и обладание-телом тесно связаны, то столкнемся с различием, относящимся к сфере нашего тела и не подрывающем его, а составляющем само его бытие. Наш телесный опыт тогда должен намного превосходить опыт тела. Опыт нашего тела предполагал бы телесность опыта, подобно тому, как наш опыт времени предполагает его темпоральность, как утверждает Гуссерль. Ницше, предшественник феноменологии тела, превозносил его как ''великий разум'' или как ''я'', управляющее им, как ''могущественный государь'' и ''неизвестный мудрец'' (Nietzsche 1980, 39-40). Но если мы продолжим развивать это видение, то можем оказаться в лабиринте тела только чтобы обнаружить, как сила этого повелителя влияет на различные части тела и как работает неведомая его мудрость. В \textit{Philosophie am Leitfaden des Leibes} Ницше старые слова ''объект'', ''субъект'', ''интерсубъективность'', логос и пафос были представлены по-новому. Далее я хотел бы развить именно это видение, опираясь на три ключевые феноменологические темы --- интенциональность, самосознание и интерсубъективность --- чтобы показать, как традиционный язык сознания превращается в новый язык телесного опыта.

\section{Смысл и аффекция}

Ранее, когда мы обсуждали общую сферу смысла, мы положили интенциональность основой феноменологии, как Фрейд положил бессознательное основой психоанализа. Но в обоих случаях ключевое слово --- слово далеко не последнее: оно не решает задачу, напротив --- открывает путь к огромному количеству вопросов. Если мы предположим, что нечто всегда дается как нечто, то есть схватывается, понимается или интерпретируется в определенном смысле и определенным образом, то возникает вопрос: как во всем этом фигурирует тело?

Конечно, наше тело участвует в работе опыта многими способами, но, похоже, оно не автор всех преднамеренных действий, приписываемых сознанию, включая сознание нашего тела. Именно Платон первым настаивал, что видят не глаза, а душа посредством глаз. Но выбор между душой и глазами может оказаться ложной дихотомией. Обратимся снова к формуле ''нечто как нечто''. Как насчет нечта, воспринимающегося как нечто? Изначально Гуссерль использует аристотелевские термины, говоря о некоем \textit{hyle} --- том, из чего нечто образуется в результате намеренных действий. Это может быть верно на уровне нормальности, где вещи являются такими, какие они есть и как таковые известны. Здесь мы можем отличить повторяемую форму от различных материалов, в которых она реализуется, как это делает Аристотель, говоря о столах или домах и их материалах. Но несколько иначе все выглядит, когда мы добираемся до истоков этих переживаний, где вещи становятся только тем, чем они должны быть. Взять, например, восприятие: оно начинается не с акта наблюдения --- напротив, оно возникает в результате привлечения внимания из-за чего-то, что меня поражает. Или действия: они возникают в ситуациях, отмеченных как нечто привлекающее или отталкивающее, пугающее или искушающее нас. Наши действия, как уже признает Гуссерль, скорее инсценированы, чем сделаны на самом деле (см. \textit{Hua IV: 98, 259, 336}). Они проходят через фазы колебаний и репетиций, прежде чем приобретут четкую форму. И память функционирует аналогично: как показывает Ницше в своей \textit{Генеалогии морали}, мы помним только о том, что причиняет нам боль, и наша спонтанная память, тщательно описанная Бергсоном и Прустом, следует нашим желаниям, а не нашей воле. Наше прошлое вновь и вновь одолевает нас, прежде чем произойдет процесс явного осознания [\textit{Vergegenwärtigung}]. Наконец, даже наше мышление начинается с идей, которые просто приходят нам в голову, с \textit{Einfälle}, или, в терминологии Лакана, \textit{incidences}. Навязчивые идеи и бред --- не просто искажения свободной, самодавлеющей мысли.

Глубины и фон претерпевания, подробно обсуждавшиеся ранее, приобретают цвет и весомость только если рассматривать телесные компоненты чувственного опыта вне превращений их в блеклую копию ментальных действий и предрасположенностей. Все, что является как нечто, должно быть описано не просто как воспринимающее смысл и обладающее им, но как смысл провоцирующее и само по себе не осмысленное --- то есть как нечто, что нас трогает, воздействует на нас, стимулирует, удивляет, преступает. Я называю это пафосом, \textit{Widerfahrnis} или аф-фектом, где дефис указывает, что с нами происходит нечто нами не инициированное. В немецком языке приставка \textit{an-} используется для обозначения ряда глаголов, таких как \textit{Angehen, Anblicken, Anreden} или \textit{Antun}, и более поздний Гуссерль следовал этим лингвистическим маркерам при разработке своей теории аффекции. Можно, конечно, спросить, достаточно ли он продвинулся по этому пути, чтобы сделать надлежащие выводы. Но в любом случае, за преднамеренными действиями, приписываемыми субъекту с точки зрения авторства или происхождения, скрываются события нас одолевающие, с нами случающиеся. Они не относятся ни к описанию от первого лица как субъективные, совершаемые нами, действия, ни к описанию от третьего лица как объективный, регистрируемый или осуществляемый извне процесс. Они требуют языка, в котором ид или \textit{Es} переплетены со мной [\textit{mich}] или привязаны ко мне [\textit{mir}]. В этом смысле эго, появляющееся в винительном или дательном падеже, предшествует эго в падеже именительном. С самого начала я вовлечен, но не в качестве ответственного автора или агента. Чтобы подчеркнуть пассивный статус так называемого субъекта, я использую термин ''пациент'', который следует понимать буквально. Соответствующим статусом был бы статус отвечающего, который от-вечает на то, что его или ее поражает. \textit{Wodurch}, то есть то, что на нас воздействует, предстает как таковое только в терминах \textit{Worauf}, то есть чего-то, на что мы отвечаем. \textit{Nachträglichkeit}, то есть последующий характер ответа, который нам предлагается дать, соответствует \textit{Vorgängigkeit}, то есть предшествующему одолевающему нас пафосу. Таким образом, во всем, что мы делаем, нас отделяет от нашего собственного истока подлинный темпоральный лаг. Без такой радикальной формы темпоральности, соединяющей то, что разделяет и разъединяющей то, что соединяет, мотив пассивности потерял бы свою силу. Он оказался бы либо внутри, как собственный предел субъекта --- как признак его конечности, например, либо снаружи, как грубый остаток реальности --- как некодированный стимул, например.

На данный момент мы уже частично ответили на вопрос о теле и его роли. Подверженность воздействию того, что мне чуждо [\textit{Ichfremdes}] не зависит ни от моего знания, ни от моего желания (то есть от моего сознания), она указывает на мое тело. Сфера телесности включает все, что, несомненно, имеет отношение ко мне, но не совершается мной. В этом смысле тело, будучи названо Мерло-Понти в \textit{Феноменологии Восприятия} ''врожденным комплексом'' (Merleau-Ponty 1945, 99), приближается к бессознательному, к ''эго тела'' (\textit{Körper-Ich}) и к телесному языку симптомов из психоанализа Фрейда. Но попытка пересечь границу между феноменологией и психоанализом --- далеко идущее начинание, за которое брались многие феноменологи. Возможно, в контексте этого обсуждения будет достаточно сосредоточиться на том, что можно назвать рождением смысла из пафоса. Как и любые другие роды, это происходит через мучительный труд. В целом, наше тело кажется более человечным, чем разум, находящийся лишь в позиции функционирования. Потому мы должны спрашивать не только что компьютер может или не может делать, но и что может и не может с ним происходить, на что он отвечает, а на что не отвечает.

\section{Расщепленная самость}

Интенциональность управляет нашим взаимодействием с миром, поиск же само-сознания или само-осознания ставит нас лицом к лицу с самими собой. Встает вопрос: как выглядит эта само-референтность с точки зрения тела. В отличие от более ранних форм само-референтности --- таких как забота о себе, о которой говорили и Хайдеггер, и Фуко --- в современности доступ к себе по большей степени связывается с так называемой рефлексией. Однако при использовании этого заимствованного из оптики понятия возникают коварные ловушки. Попытка Локка рассматривать самость как нечто вне вещей --- и так дополнить внешний опыт внутренним --- обречена на провал, как только мы вступаем с сферу интенциональности: если я являюсь и тем, благодаря чьим намеренным действиям вещи конституируются как таковые, и тем, кому они даны, то эго сводится к лишь ''символическому концу мира'' [\textit{Endchen der Welt}] (\textit{Hua} I:63). Потому для Гуссерля рефлексия означает не какой-то внутренний опыт, а трансцендентальную рефлексию, то есть тематизацию соотношения между смыслом и действием, включая те горизонты, в которых разворачивается моя жизнь в мире. Иначе говоря, рефлексия преобразует функционирующую интенциональность --- в явную.

Эта рефлексия, разумеется, отображает многие телесные фспекты. Мы видим, слышим, осязаем и двигаем себя, наша собственная телесность так ''отсылает к себе'' (\textit{Hua} I:128). Любая само-референтная деятельность  включает различные формы само-аффекции. Когда мы смотрим на себя в зеркало, слышим свой голос или прикасаемся к острому ножу, мы удивляем сами себя. Мы захвачены собственным образом, сбиты с толку собственным голосом или врезаемся в собственную плоть. Но пока мы интерпретируем это ''бытие для себя'' как ''интенциональное отношение к себе'' (\textit{Hua} I:81), мы запутываемся в бесконечных итерациях: переходим от одной рефлексии к другой, так и не достигая объекта рефлексии. Из этого сизифова труда можно извлечь кое-что важное: телесная самость --- не-вещь [\textit{Nicht-Ding}], никогда не ''телесно присутствующая'' [\textit{leibhaft gegenwärtig}] подобно вещам. Цепь рефлексий опирается на дорефлексивное ядро живого настоящего, могущее быть схвачено только впоследствии, посредством своего рода ''пост-осознанности'' [\textit{Nachgewahren}] (\textit{Hua} VIII:89). Многолетние дискуссии вокруг самосознания сполна продемонстрировали, что каждая попытка построить самосознание на рефлексии предполагает наличие самости, ради объяснения которой и затевалась. Этой дилеммы не избегает даже так называемый \textit{авто-поэзис} --- самопостроение, если понимать буквально.

Но есть и другие --- гораздо теснее связанные с проблемой тела --- проблемы. Конечно, мы знакомы с многчисленными процессами и состояниями: такими как кровяное давление, гормональный баланс, возбуждение нейронов и, наконец, функционирование ''моего мозга'', центральной части ''моего тела''. Все это относится к моему телу, однако не участвует в вышеупомянутых формах чувственной рефлексии. Нет смысла заявлять, будто наш мозг размышляет о себе как если бы был потомком древнего гомункула. Утверждая, что мозг выбирает или оценивает определенные стимулы, неврологи на самом деле имеют в виду операции, не требующие оперативной самости. Рекурсивные процессы, создаваемые операционными циклами, довольно сильно отличаются от обычных рефлексивных актов. Может возникнуть соблазн противостоять этому вторжению физики в сферу нашего живого тела и свести весь нейрофизиологический аппарат к натуралистическому подходу, оставив нетронутым так называемый персоналистический подход. Но таким жестким разграничением мы не избавимся от проблемы ''душа-тело''. Конечно, мозг как таковой --- неврологическая конструкция, но не то, что следовало бы называть ''функционирующим мозгом'', проявляющим себя не просто в измеримых данных, а в сильной головной боли, например. Активный мозг принадлежит нашему живому телу и не является каким-то физическим приложением. Простое различие между двумя установками, предложенное Гуссерлем, все еще слишком тесно связано с предполагаемым сознанием нашего тела и так возвращает нас на картезианский путь. Гораздо убедительней звучит определение Гуссерлем тела как ''точки перехода'' [\textit{Umschlagstelle}] между чувственной и естественной причинностью (\textit{Hua} IV:286). Один из таких переходов мы испытываем при повседневной усталости: как утверждает Поль Валери в своей книге \textit{Cahiers I} (11:37), ''Усталость превращает 'тело' в странную вещь'' (''Par la fatigue le 'corps' devient chose étrangère''). Мое собственное тело приобретает так черты тела чуждого. Более того, технологии растущим количеством техник воздействия на тело порождают новые --- требующие настоящей феноменотехники --- проблемы. Это видно в повседневной жизни: когда я слышу свой голос из динамика или вижу свое лицо на видео, я попадаю в ситуацию, когда близость и отдаленность переплетаются, когда все прямые отражения своеобразно искажаются. Старые фотоальбомы, в которых я вижу себя ребенком, показывают, кем я был и кем мог бы стать. Мальчик подает ''надежды, еще не его собственные, неопределенные ожидания желаемого славного будущего, тянущиеся к нему распростертыми крыльями золотой сети'', и так же вопрошает Ульрих, ''Человек без свойств'': ''Что бы могло быть достаточно прочным, что могло бы удержать нас от распада на части?'' (Musil 1978, 648).

Это реальные проблемы, и соблазн их отбросить --- найти прибежище в чистой самости, в чистом ощущении своего тела, имманентности жизни, возвышенной и всему выталкивающему нас за наши пределы (будь то водоворот времени, экстаз интенциональности или вызовы, бросаемые нам другими) предшествующей само-аффекции --- телесный нарциссизм, встречающийся у Мишеля Генри или Германа Шмитца, основанный на несостоятельных предпосылках. Восприятие тела как такового всегда предполагает определенную дистанцию к своему телу, которую Плеснер называет ''эксцентричной позиционностью''. Чистое тело --- \textit{Leib ohne Körper} --- подчиняется картезианскому правилу \textit{sentio, ergo sum}, но осознать свои чувства и дать им какое-то выражение --- нечто большее, чем просто им подчиниться.

Потому следует искать другие пути, и настаивание на расщепленности самости может быть одним из них. Декарт здесь снова поможет нам прояснить, что поставлено на карту: от его позиции нельзя просто отмахнуться, ее необходимо учитывать и трансформировать, иначе она будет возвращаться вновь и вновь, как всякое вытесненное. Согласно картезианской точке зрения, есть два фундаментальных способа бытия: мыслящий разум и протяженные вещи. В терминах различных способов восприятия это значит: умы видят, но не могут быть увидены, а вещи видимы, но не могут видеть. Тело не вписывается в эту дуалистическую схему: оно и то, и другое, видящее и видимое, слышащее и слышимое, касающееся и касаемое, движущееся и движимое. Но видящий и видимое, как \textit{cogito} и \textit{cogitatum}, никогда не совпадают. Несовпадение это характеризует само бытие нашего тела, к себе обращающегося и в то же время себя избегающего. Первоначальный пример чувственной рефлексии не следует рассматривать как эстетическую версию аристотелевского νόησις νοήσεως, конечным результатом которого становится превращение мышления о мышлении в видение видения, όρασις όράσεως, как у Плотина (\textit{En.} 5.3.8). От глаз ускользает не видение, а событие становления видимым: что-то поражает меня и привлекает мое внимание, и ускользает здесь мой взгляд, а не глаза (которые я, разумеется, могу наблюдать в зеркале и форму и цвет которых могу всегда определить). Взгляд, отвечающий на возникающее, не имеет цвета и в этом смысле невидим. Замечая, что в зеркале мое лицо отражено, я думаю о себе как о смотрящем на себя из зеркала --- а не как о лишь отражении --- потому что отражение находится на своем месте. Но сравнивая это ускользание со слепым пятном на сетчатке, следует быть осторожным и не перепутать невидимое изменение видимого с лишь пробелом в поле зрения, как это показывает Эрнст Мах в своем эскизном автопортрете.

Далее, наше тело становится не только видимым, но и осязаемым, движимым, выразительным и либидозным. Рваные линии моего тела расходятся в разные стороны. Юный Нарцисс влюбляется в свое отражение и падает насмерть. Его преследует нимфа Эхо, которая отвечает лишь повторением сказанного. Это дважды неудачное свидание, которое Овидий представляет нам в мифической форме, учит многому о телесности. Эхо возникает как акустическая версия эффекта зеркала: мы слышим, как говорим, хотя говорим именно мы. Здесь любопытен феномен аутоэхолалии --- когда младенцы пытаются имитировать собственные вокализации,--- а также особый случай эхолалии, повторение услышанного, с помощью которого пожилые люди пытаются компенсировать свою ментальную регрессию. Вербальные галлюцинации, при которых пациенты слышат голоса как бы извне (хотя сами их издают), совершенно непонятны, если понимать голос как всегда нам принадлежащее. Картезианцам следует попытаться объяснить возможность галлюцинаций, иначе психически больные низводятся ими до уровня странных живых существ, и психиатрия теряет --- как это уже часто бывало --- свой человеческий облик.

Что же касается движущегося тела, то оно играет роль \textit{basso continuo} для всех видов действий. Это старая проблема. Согласно Аристотелю, тело приводится в движение целями, притягивающими или отталкивающими наши стремления. Когда разум отделен от эмоций, он ничего не движет (см. \textit{Nic.Ethics} 6.2.1139a35-36). Движения и эмоции идут рука об руку. В наше время телесные движения также рассматриваются дуалистически. В своем трактате \textit{Страсти души} Декарт проводит четкое различие между действием --- когда я изнутри двигаю нечто,--- и страстью --- когда нечто меня движит извне. Он рассматривает и интересные предельные случаи, такие как неугомонность или беспокойство, которые не привносятся объектами извне, а пробуждаются (\textit{Passiones animae 2.51}). Ключевая альтернатива здесь следующая: двигать или быть движимым, или, в Кантовских терминах, каузальность свободы или каузальность природы. Но если практическое поведение --- как и прочие его формы --- начинается с претерпевания и продолжается отвечанием, то мы сталкиваемся с загадкой существа, которое двигает, но не полностью самостоятельно. Это существо тогда участвовало бы и в текущем движении, вмешивалось в него, движение это также предшествовало бы нашей инициативе. Пальцы музыканта куда быстрее и чувствительнее, чем это было бы возможно при рациональном управлении ими, ноги футболиста следуют за мячом, будто руководствуясь навигатором внутри его тела. Нас увлекают (\textit{mitgerissen, comportés}) наши собственные слова и поступки и поступки других людей, так что мы не сводимся ни к движимым объектам, ни к движущим субъектам, и потому в нашу повседневную жизнь и проникает то, что Фрейд называет \textit{Fehlleistungen} --- своего рода парапраксис, который относится к нашей практике, но не сводится к добровольным решениям. Движимость без движения и движение без движимости --- всего лишь предельные случаи, в них мы либо опускаемся до состояния безжизненных вещей, либо поднимаемся в высшему, лишенному желаний состоянию богподобного Аристотелевского ''перводвигателя'' --- движущего через любовь к себе (κινεϊώς έρώμενον), но не любя (Aristotle 1984 1072b3).

Можно продолжить и показать, что все наше поведение проистекает из своего рода само-аффекции, претерпеваемой нами в ответе на нее же. Мы старше самих себя. В результате нашего рождения и череды перерождений мы не способны собрать свою \textit{Selbstvorsprung}, свое самопредшествование. Рождение смысла из пафоса, упомянутое выше, подкрепляется рождением меня самого из пафоса: мое прошлое порождает непреодолимую чуждость, экстатическую чуждость. Я выхожу за пределы себя не по воле случая, болезни или слабости, а через бытие тем, кто я есть. Эта инаковость --- результат нарушенной самореференции, включающей определенные формы самоустранения, ускользания. Связанный с собой и в то же время от себя отрезанный, я не одно и не два, но два в одном и одно в двух. Внутреннее напряжение между обоими полюсами оставляет место для крайних форм как слияния, так и расчленения. И монизм, и дуализм находят здесь определенную истинность.

Что же касается старой проблемы соотношения души и тела, то можно было бы говорить о \textit{Leibkörper} в точном смысле этого слова, как это делали Гуссерль, Шелер и Плеснер. Это сложное существо включает не только живое тело --- с помощью которого мы воспринимаем вещи и манипулируем ими, выражаем себя и кооперируемся,--- но и весь физиологический аппарат, неврологические и генетические процессы, которыми наше собственное поведение не только реализуется, но и, в некоторой степени, формируется. Все это принадлежит нам, но с точки зрения уменьшающейся близости и увеличивающегося расстояния. Я могу говорить об определенном мозге как о принадлежащем мне, нужно только учесть, что принадлежать мне --- не \textit{eo ipso} значит быть в моем распоряжении. Мы направляем свой взгляд, протягиваем руки или ускоряем шаг, но не задерживаем дыхание и не повышаем кровяное давление. Это не исключает того факта, что так называемые высшие формы мышления и волеизъявления по-разному взаимодействуют с физиологическими процессами. Сильные аффекты, такие как удивление или тревога, врываясь в нашу обыденность и провоцируя ответы, как правило окружены особой аурой телесности. Это верно и для θαυμάζειν в \textit{Теэтете} Платона (155c), сопровождающегося головокружением, и для \textit{admiratio} в \textit{Страстях души} Декарта (2.74), от которого будто каменеем. Если не подчинять мышление и речь исключительно стандартам ясности или корректности и делать акцент на глубоких убеждениях и устоявшихся привычках, то ничто не помешает признать воплощенную форму мышления, \textit{leibhaftig}. Невозможно мыслить и желать без глубоких убеждений и привычек, и они далеко не всегда в нашем свободном распоряжении. Только такой писатель, как Достоевский, мог сделать смелое предположение, что в темноте ночи мы верим иначе, чем при свете дня, и именно Ницше обвинил философов в выращивании идей сидя [\textit{ersessen}] вместо собирания их во время прогулок [\textit{ergangen}]. По-немецки мы говорим о \textit{Gedankengänge}, или ''тропах мысли''. \textit{Ambulo ergo sum} Гассенди следует понимать не как обратную сторону \textit{cogito ergo sum} Декарта, а как его воплощение.

Непрерывная шкала близости и отдаленности оставляет место для самых разнообразных форм патологии. Позвольте здесь упомянуть случаи деперсонализации, когда рука пациента лежит на столе и видится им подобно камню, случаи шизофрении, когда пациент отрезан от собственных идей, случаи травмы, когда он зациклен на произошедшем с ним, на что не в состоянии адекватно отреагировать. Такие патологические трещины не выходят полностью за рамки нормального расщепления нашей самости, всегда пронизанной и одержимой определенными аномалиями. К подобным патологиям может привести попытка избежать хаоса чрезмерной нормализацией. С точки зрения тела, болезнь требует различных терапевтических вмешательств, направленных на преодоление картезианского разделения на соматические и психические заболевания. Следует проводить другие различия, например между формами болезней на переферии нашей телесной самости и формами другими, затрагивающими суть нашего существования. Называя человека ''мыслящим тростником'', Паскаль был ближе к истине о теле, чем Декарт, который лишь позднее, перейдя от порядка разума к порядку жизни, признал смешение души и тела.

\section{Другой как мой двойник}

Чуждость во мне открывает путь к чуждости Другого и не позволяет вернуться на путь традиционного дуализма, согласно которому некоторые тела следует называть телами других только в отношении других разумов.\footnote{Выражение ''другие разумы'' мы встречаем в англосаксонской литературе, а именно в названии статьи Дж.Л.Остина, опубликованной в 1946г.} Всякий раз, когда я начинаю с сознания чуждого (\textit{Fremdbewußtsein}), я конституирую чуждость Другого конституируется на основе моей ''сферы владений'', ''внутри и посредством того, что мое'' (\textit{Hua} I:131). Чуждость Другого всегда будет производной от своей: Другой предстает как альтер-эго, как второй я. Но как только мы становимся на точку зрения тела и исходим из телесной самости, которая ''в собственном доме не хозяин'', Другой предстает как со-изначальный мне и в какой-то степени как мне предшествующий. Чтобы избежать традиционного --- основанного на субъективизме, лежащем в основе всего и каждого, с кем мы сталкиваемся --- эгоцентризма, нужно последовать примеру Мерло-Понти, заменившего термин ''интерсубъективность'' на ''интеркорпоральность''. Изменение базовой схемы меняет и характер межлежащего: диа- и взаимо- в диа-логе, \textit{entre-tien} или взаимо-отношении.

Рассматривая статус Другого, мы сталкиваемся с двумя особенно важными аспектами. Во-первых, Другой изначально возникает ни как что-то --- которое я затем преобразую в кого-то через эмпатию или аналогию,--- ни как кто-то --- намерения которого я понимаю, интерпретирую и разделяю. Инаковость и чуждость Другого заявляет о себе в терминах пафоса, как специфическая \textit{Fremdaffektion}. Окружающие врезаются в наше прежде чем мы сможем спросить, кто они и что значат их выражения. Чуждость Другого одолевает и удивляет нас, преступая наши намерения прежде чем будет так или иначе понята. Второй аспект проистекает из того факта, что на меня воздействует не только другое эго или другой субъект --- кто-то, отличный от меня --- но и \textit{meinesgleichen}, \textit{mon semblable}, похожий на меня и в то же время несравнимый, \textit{unvergleichlich, hors de série}. Появление Другого не следует понимать как нечто само собой разумеющееся, как это делают многие теоретики языка, общества и культуры. Есть не только ''чудо разума'', о котором говорил уже Лейбниц и которое позже было подхвачено Гуссерлем и Мерло-Понти, но и ''чудо Другого'' в смысле отсутствия для него достаточного основания. У нас есть яблоки, столы и компьютеры --- в том числе тот, что сейчас передо мной,--- но ни одна из этих вещей не имеет себе равных. Бытие кем-то --- а не просто чем-то --- подразумевает наличие равных, но эти равные остаются неравными из-за своей уникальности, коренящейся в их воплощенном существе, в их бытии здесь и сейчас. Этот странный, невыводимый ни из чего другого, факт я называю удвоением себя в Другом и Другим. Таким образом, моя собственная экстатическая чуждость усиливается дублирующей чуждостью Другого. Рядом здесь тревожный феномен двойника, доппельгангера, исключительно тесно связанный Полем Валери с существованием Другого. В \textit{Cahiers} он пишет: ''Другой --- еще один подобный мне, или, может, мой двойник --- самая притягательная бездна, самый оживляющий вопрос, самое злостное препятствие, что одно не дает всему остальному запутаться, разорваться'' (Valéry 1973, 499).

Этот странный вид чуждости можно проиллюстрировать теми же примерами, которые использовались ранее в объяснении собственной чуждости. Чуждый взгляд, которому я открыт, заключается в чувстве увиденности прежде собственного видения Другого как человека, видящего вещи, и среди них меня. Крайности это достигает в особой форме паранойи --- ощущении, что за мной постоянно наблюдают [\textit{Beobachtungswahn}]. Это нельзя свести к простому эффекту ''я вижу то, чего ты не видишь'', считаемого достаточным системными теоретиками. ''Слепое пятно'' присуще любому чуждому опыту и выходит за пределы возможностей самореферентных систем, потому что они неспособны интегрировать собствевнное функционирование. Бессмысленно утверждать, что машина чувствует себя преследуемой сенсорами других машин: человеческий другой не только ставит нас перед лицом пределов наших возможностей --- аффекция другого преступает нас как \textit{wirkende Wirklichkeit}, как аффективная реальность, предшествующая условиям возможности, анализируемая трансцендентальной философией и конкретизируемая в терминах нормативных систем или кодексов. Чуждый опыт, строго говоря, невозможен: он не делается возможным ни по моей собственной инициативе, ни в соответствии с общими правилами. Следуя Шеллингу, его можно назвать незапамятным [\textit{unvordenklich}].

То, что верно для чуждого взгляда, верно и для чуждого голоса. Чувство, что ко мне обращаются, не ограничивается слышимостью, как кто-то произносит определенные слова или издает определенные звуки. Слушать обращение к себе --- нечто большее, чем слышать, как что-то сказано Другим. Множество лингвистов и некоторые ориентированные на лингвистику философы забывают, что взаимопонимание и намерение невозможно без внимания (\textit{no in-tention without at-tention}). То, что я здесь называю пафосом или аффекцией, отмечено чертами чуждости [\textit{Ichfremdheit}] еще до того, как будет приписано кому-то предположительно его вызвавшему. Как уже было показано, взгляд и оклик не ограничиваются тем, что я становлюсь видимым и слышащим --- они включают меня как того, кого видят и к кому обращаются. Это значит, что я воспринимаю себя со стороны. Все это относится к повседневным ситуациям, в фоне которых лежит не-повседневное, более того --- это состояние глубоко укоренилось в нашем собственном происхождении, что видно, например, в процессе присвоения имен. Так называемое имя собственное --- имя на самом деле наполовину чуждое, потому что дается другими. Имя, которым меня называют --- мое, потому что я откликаюсь на него, в то время как от объекта подобного отклика, и тем самым принятия ярлыка, не ожидается. Более того, в моем обществе сформировалось множество привычек и манер поведения, и даже мой родной язык заимствован у других людей --- я знаю его буквально понаслышке. Наше тело сексуализировано и в значительной степени сформировано желанием Другого, и желание здесь следует понимать двояко: как исходящее от Другого и как направленное к Другому. Учитывая, что Другой как бы имплантируется в нас из раннего симбиоза ребенка с матерью, можно последовать за психоаналитиком Жаном Лапланшем и заявить о главенстве соблазнительного другого (Laplanche 1997). Нельзя достичь Другого qua Другого, не начиная с него или нее.

Впрочем, из этого не следует, что нужно обратить вспять унаследованный эгоцентризм и заменить его аллоцентризмом, ведь и с Другого мы не можем начать, не обращаясь неявно к себе. Чуждая любовь не сводится к само-любви, хотя и не мыслится без нее. Более того, подобно тому, как мое тело как Leibkörper занимает непрерывную шкалу близости и дистанции, мое социальное тело принадлежит мне лишь в большей или меньшей степени, зависимо от меняющейся степени близости и дистанцированности по отношению к другим. Интеркорпоральность подразумевает переплетенность своего и чуждого, включенность каждого в нексус, в \textit{Geflecht}, как утверждали и Норберт Элиас, и Мерло-Понти. В случае с Мерло-Понти можно даже говорить о социальном синкретизме (см. \textit{Antwortregister}, часть III, глава 8). Нет никаких заранее готовых индивидов --- только процесссы индивидуации, придающие нашей телесной самости определенную степень анонимности и наделяющие ее телесной типичностью. То, что мы чувствуем, воспринимаем, делаем или говорим, переплетается подобно тибетскому ковру с тем, что чувствуют, воспринимают, делают или говорят другие. То, что делает или говорит один [\textit{Man}], не обязательно опускается до уровня неаутентичности или вульгарности [\textit{Uneigentlichkeit} или \textit{Vulgarität}] --- оно составляет общий фон того, что мы делаем или говорим от своего имени, и становится неаутентичным только если человек ведет себя таким образом, что отождествляет себя с общественным мнением или общепринятой моралью. И это еще не все. Тот, кто получил свое имя от других и, значит, им --- в отличие от родимого пятна --- не владеет, сохраняет безымянность в самой глубине своей души. Древние табу на использование имен или создание изображений указывают на нашу неприкосновенность, предполагающую чуждость и отстраненность тела. Следует, значит, исходить из функционирования нашего тела как ''переходной точки'' [\textit{Umschlagstelle}], и не только потому, что собственное обращается чуждым и наоборот: говорить о ''моем собственном теле'' или о \textit{corps propre} --- значит говорить лишь половину правды. В своей теории романа Михаил Бахтин развивает идею ''внутреннего диалогизма'' речи и утверждает, что каждое слово --- ''наполовину чуждое слово'', потому что ''заряжено --- даже перегружено --- чуждыми намерениями'' (Бахтин, 1979, с. 185). Аналогично и мое собственное тело можно описать как наполовину чуждое --- заряженное исходящими от других, чуждыми, намерениями, желаниями, проекциями, привычками, аффекциями и преступлениями.

В конце концов, наш телесный и воплощенный опыт показывает, что я ощущаю Другого в себе и себя в Другом еще до того, как мы встретимся друг с другом. Как подчеркивает Мерло-Понти, другие возникают во мне и рядом со мной прежде встречи со мной лицом к лицу. В \textit{Прозе мира} он пишет: ''Другой . . . потому всегда на границе мной видимого и слышымого, он по эту сторону от меня, рядом со мной или позади, но не в том месте, которое мой взгляд выглаживает и опустошает от всякого 'внутреннего'. . . . Он будто иногда ощущаемый больным рядом двойник, похожий на него как на брата, на которого он не мог бы положиться, не заставив исчезнуть'' (1969, 186; 1973, 134). И даже когда я смотрю на Другого, я не нахожусь на одном уровне с ним. Обмен происходит по наклонной плоскости, и верхнюю часть всегда занимает Другой. Согласно Левинасу, зов чуждого исходит с ''высоты'', поскольку превосходит наши ожидания, желания и идеи и преступает нас прежде, чем мы успеваем занять надлежащую позицию. Впрочем, здесь следует остерегаться недоразумений.

\end{document}
